\subsection{Measure induced on a locally compact subspace}

Let X be a locally compact space, \(\mu\) a measure on X, and Y a locally compact subspace of X. Since Y is the intersection of an open set and a closed set in X (GT, I, §9, No.~7, Prop. 12), it is \(\mu\) -measurable (No.~1, Cor. of Prop. 3). For every function \(g \in \mathcal{K}(\mathrm{Y};\mathbf{C})\) let \(g'\) be the function, defined on all of X, equal to g on Y and to 0 on X - Y; let us show that \(g'\) is \(\mu\) -integrable. We can restrict ourselves to the case that g is real and \(\geqslant 0\) (on writing g as a linear combination of such functions); since \(g'\) is bounded and has compact support, it suffices to show that \(g'\) is \(\mu\) -measurable (No.~6, Th. 5); but this follows from the fact that \(g'\) is upper semi-continuous on X (No.~5, Cor. of Prop. 8). We may therefore make the following definition:

DEFINITION 4. --- Given a locally compact subspace Y of a locally compact space X, we call measure induced on Y by a measure \(\mu\) on X, and denote by \(\mu_{Y}\) or \(\mu|Y\) , the measure defined by the formula

\[
\int g d \mu_ {\mathrm{Y}} = \int g ^ {\prime} d \mu\tag{1}
\]

for every function \(g \in \mathcal{K}(\mathrm{Y};\mathbf{C})\) , where \(g'\) denotes the function equal to \(g\) on \(\mathrm{Y}\) and to 0 on \(\mathrm{X} - \mathrm{Y}\) .

Example. --- Let \(\mu\) be Lebesgue measure on R, I any interval in R; I is a locally compact subspace of R, and the measure induced by \(\mu\) on I is the linear form

\[
g \mapsto \int_ {a} ^ {b} g (x) d x
\]

on \(\mathcal{K}(\mathrm{I};\mathbf{C})\) , where \(a\) and \(b\) are the endpoints (finite or not) of I (cf.~§ 4, No.~4, Example), in other words, what we have called Lebesgue measure on I.

When Y is an open subspace of X, Def. 4 coincides with the definition of the measure induced by \(\mu\) on Y (or the restriction of \(\mu\) to Y) given in Ch. III, §2, No.~1: indeed, for every function \(g \in \mathcal{K}(\mathrm{Y}; \mathbf{C})\) the function \(g'\) is then continuous on X.

We shall study integration with respect to an induced measure in detail in Ch. V, §7, and until then we shall need only the following results:

Lemma 2. --- Let \(\mu\) be a positive measure on \(X\) , and let \(K\) be a compact subset of \(X\) .

\begin{enumerate}
\def\labelenumi{(\roman{enumi})}
\item
  For every compact (resp. open) subset H of K, \(\mu_{\mathrm{K}}(\mathrm{H}) = \mu (\mathrm{H})\)
\item
  For a subset \(\mathbf{N}\) of \(\mathbf{K}\) to be \(\mu_{\mathbf{K}}\) -negligible, it is necessary and sufficient that it be \(\mu\) -negligible.
\item
  If S is the support of \(\mu_{\mathrm{K}}\) , then \(\operatorname{Supp}(\mu_{\mathrm{S}}) = \mathrm{S}\) .
\setcounter{enumi}{0}
\item
  We can restrict ourselves to the case that H is compact. Denote by f the characteristic function of H in the space K; f is upper semi-continuous, hence is the lower envelope of a decreasing directed family \((g_{\alpha})\) of functions in \(\mathcal{K}_{+}(\mathrm{K})\) ; we have \(\mu_{\mathrm{K}}(\mathrm{H}) = \inf_{\alpha} \int g_{\alpha} d\mu_{\mathrm{K}}\) (§4, No.~4, Cor. 2 of Prop. 5). If \(g_{\alpha}^{\prime}\) is the function equal to \(g_{\alpha}\) on K and to 0 on X - K, then \(g_{\alpha}^{\prime}\) is upper semi-continuous, and the lower envelope of the decreasing directed family \((g_{\alpha}^{\prime})\) is the characteristic function \(\varphi_{H}\) of H in the space X; therefore
\end{enumerate}

\[
\mu (\mathrm{H}) = \inf _ {\alpha} \int g _ {\alpha} ^ {\prime} d \mu = \inf _ {\alpha} \int g _ {\alpha} d \mu_ {\mathrm{K}} = \mu_ {\mathrm{K}} (\mathrm{H})
\]

by (1).

\begin{enumerate}
\def\labelenumi{(\roman{enumi})}
\setcounter{enumi}{1}
\item
  If N is \(\mu\) -negligible then, for every \(\varepsilon > 0\) , there exists a relatively compact open neighborhood U of N in X such that \(\mu(\mathrm{U}) \leqslant \varepsilon\) ; since \(\mathrm{K} - (\mathrm{U} \cap \mathrm{K})\) is compact, it follows from (i) that \(\mu_{\mathrm{K}}(\mathrm{U} \cap \mathrm{K}) \leqslant \mu(\mathrm{U}) \leqslant \varepsilon\) , thus N is \(\mu_{K}\) -negligible. Conversely, if N is \(\mu_{K}\) -negligible then there exists an open neighborhood V of N in X such that \(\mu_{\mathrm{K}}(\mathrm{V} \cap \mathrm{K}) \leqslant \varepsilon\) ; by (i), we have \(\mu(\mathrm{V} \cap \mathrm{K}) = \mu_{\mathrm{K}}(\mathrm{V} \cap \mathrm{K})\) , therefore \(\mu(\mathrm{N}) = 0\) since \(\varepsilon\) is arbitrary.
\item
  For every open set U in K that intersects S, by hypothesis \(\mu_{\mathrm{K}}(\mathrm{U}\cap\mathrm{S})\neq0\) , therefore \(\mu(\mathrm{U}\cap\mathrm{S})\neq0\) by (i), and since \(U\cap S\) is open in S, \(\mu_{\mathrm{S}}(\mathrm{U}\cap\mathrm{S})\neq0\) by (i); since every nonempty open set in S is of the form \(U\cap S\) , where U is open in K, this proves that \(\operatorname{Supp}(\mu_{\mathrm{S}})=\mathrm{S}\) .
\end{enumerate}

Lemma 3. --- Let Y be a locally compact subspace of X; for every measure \(\mu\) on X, \(|\mu_{\mathrm{Y}}| = |\mu|_{\mathrm{Y}}\) .

Let \(f\) be a function in \(\mathcal{K}_{+}(\mathrm{Y})\) and \(\varepsilon\) any number \(>0\) ; by definition, there exists a function \(g \in \mathcal{K}(\mathrm{Y};\mathbf{C})\) such that \(|g| \leqslant f\) and \(|\mu_{\mathrm{Y}}|(f) \leqslant |\mu_{\mathrm{Y}}(g)| + \varepsilon\) . Denoting by \(f'\) and \(g'\) the functions obtained by extending \(f\) and \(g\) , respectively, to be 0 on \(\mathrm{X} - \mathrm{Y}\) , we have \(\mu_{\mathrm{Y}}(g) = \mu(g')\) and, since \(|g'| \leqslant f'\) ,

\[
\left| \mu \left(g ^ {\prime}\right) \right| \leqslant | \mu | \left(\left| g ^ {\prime} \right|\right) \leqslant | \mu | \left(f ^ {\prime}\right) = | \mu | _ {Y} (f),
\]

whence \(|\mu_{\mathrm{Y}}|(f)\leqslant |\mu |_{\mathrm{Y}}(f) + \varepsilon\) and, since \(\varepsilon\) is arbitrary,

\[
| \mu_ {\mathrm{Y}} | (f) \leqslant | \mu | _ {\mathrm{Y}} (f).
\]

On the other hand, let K be the support of f and let U be a compact neighborhood of K in X such that \(|\mu|(\mathrm{U}-\mathrm{K}) \leqslant \varepsilon\) ; by Urysohn's theorem, there exists a function \(f_{1} \in \mathcal{K}_{+}(\mathrm{X})\) , extending f, with support contained in U and such that \(\|f_{1}\| = \|f\|\) . There exists a function \(h_{1} \in \mathcal{K}(\mathrm{X}; \mathbf{C})\) such that \(|h_{1}| \leqslant f_{1}\) and \(|\mu|(f_{1}) \leqslant |\mu(h_{1})| + \varepsilon\) . If h is the restriction of \(h_{1}\) to Y, then \(h \in \mathcal{K}(\mathrm{Y}; \mathbf{C})\) , \(|h| \leqslant f\) and \(\mu(h_{1}) - \mu_{\mathrm{Y}}(h) = \mu(h_{1}\varphi_{\mathrm{U}-\mathrm{K}})\) , therefore

\[
\left| \mu \left(h _ {1}\right) - \mu_ {\mathrm{Y}} (h) \right| \leqslant \| f \| \cdot | \mu | (\mathrm{U} - \mathrm{K}) \leqslant \varepsilon \| f \|;
\]

moreover, we have similarly

\[
| \mu | (f _ {1}) - | \mu | _ {\mathrm{Y}} (f) = | \mu | (f _ {1} \varphi_ {\mathrm{U} - \mathrm{K}}) \quad \text { and } \quad \left| | \mu | (f _ {1}) - | \mu | _ {\mathrm{Y}} (f) \right| \leqslant \varepsilon \| f \|.
\]

From this we infer that

\[
| \mu | _ {\mathrm{Y}} (f) \leqslant | \mu_ {\mathrm{Y}} (h) | + \varepsilon (1 + 2 \| f \|) \leqslant | \mu_ {\mathrm{Y}} | (f) + \varepsilon (1 + 2 \| f \|)
\]

and since \(\varepsilon\) is arbitrary, \(|\mu|_{\mathrm{Y}}(f) \leqslant |\mu_{\mathrm{Y}}|(f)\) , which completes the proof.

\subsection{\texorpdfstring{\(\mu\) -dense families of compact sets}{mu -dense families of compact sets}}

PROPOSITION 12. --- Let \(\mu\) be a measure on a locally compact space \(X\) , A a \(\mu\) -measurable subset of \(X\) , and \(\mathfrak{K}\) a set of compact subsets of A satisfying the following conditions:

\((\mathrm{PL}_{\mathrm{I}})\) Every closed (hence compact) subset of a set of \(\mathfrak{K}\) belongs to \(\mathfrak{K}\) . \((\mathrm{PL}_{\mathrm{II}})\) Every finite union of sets of \(\mathfrak{K}\) belongs to \(\mathfrak{K}\) .

The following four properties are then equivalent:

\begin{enumerate}
\def\labelenumi{\alph{enumi})}
\item
  For a subset B of A to be locally \(\mu\) -negligible, it is necessary and sufficient that \(|\mu|^*(\mathrm{B} \cap \mathrm{K}) = 0\) for all \(\mathrm{K} \in \mathfrak{K}\) .
\item
  For every compact subset \(K_{0}\) of A and every \(\varepsilon > 0\) , there exists a set \(K \in \mathfrak{K}\) , contained in \(K_{0}\) and such that \(|\mu|(\mathrm{K}_{0} - \mathrm{K}) \leqslant \varepsilon\) .
\item
  For every compact subset B of A, there exists a partition of B formed by a \(\mu\) -negligible set N and a sequence \((\mathrm{H}_{n})\) of compact sets belonging to \(\mathfrak{K}\).
\item
  For every compact subset B of A, there exists an increasing sequence \((\mathrm{K}_{n})\) of compact sets belonging to \(\mathfrak{K}\), contained in B and such that the set \(N = B - \bigcup K_{n}\) is \(\mu\) -negligible.
\end{enumerate}

It is immediate (No.~2, Prop. 5) that \(d\) implies \(a\) ; \(c\) implies \(d\) on taking \(\mathrm{K}_n\) to be the union of the \(\mathrm{H}_p\) for \(p \leqslant n\) and citing \((\mathrm{PL}_{\mathrm{II}})\) . To prove that \(b\) implies \(c\) , one defines recursively a sequence \((\mathrm{H}_p)\) of sets of \(\mathfrak{K}\) such that \(\mathrm{H}_{n+1} \subset \mathrm{B} - \bigcup_{p \leqslant n} \mathrm{H}_p\) and \(|\mu| (\mathrm{B} - \bigcup_{p \leqslant n} \mathrm{H}_p) \leqslant 1/n\) ( \(\S 4\) , No.~6, Th. 4).

It remains to prove that \(a)\) implies \(b)\) . Let us argue by contradiction, and suppose that the supremum \(\alpha\) of the numbers \(|\mu|(\mathrm{K})\) , where K runs over the set of subsets of \(\mathrm{K}_0\) belonging to \(\mathfrak{K}\) , is \(< |\mu|(\mathrm{K}_0)\) . By \((\mathrm{PL}_{\mathrm{II}})\) , there exists an increasing sequence \((\mathrm{L}_n)\) of compact subsets of \(\mathrm{K}_0\) , belonging to \(\mathfrak{K}\) and such that \(\sup_n |\mu|(\mathrm{L}_n) = \alpha\) . Set \(\mathrm{B} = \bigcup_n \mathrm{L}_n\) ; B is integrable and \(|\mu|(\mathrm{B}) = \alpha\) , therefore \(|\mu|(\mathrm{K}_0 - \mathrm{B}) = |\mu|(\mathrm{K}_0) - \alpha > 0\) . On the other hand, we shall see that for every set \(\mathrm{K} \in \mathfrak{K}\) , \(|\mu|(\mathrm{K} \cap (\mathrm{K}_0 - \mathrm{B})) = 0\) , which, by virtue of \(a)\) , will imply a contradiction. Indeed, if there existed a set \(\mathrm{K} \in \mathfrak{K}\) such that \(|\mu|(\mathrm{K} \cap (\mathrm{K}_0 - \mathrm{B})) > 0\) , then there would exist a compact subset H of \(\mathrm{K} \cap (\mathrm{K}_0 - \mathrm{B})\) such that \(|\mu|(\mathrm{H}) > 0\) . By \((\mathrm{PL}_{\mathrm{I}})\) , we would have \(\mathrm{H} \in \mathfrak{K}\) , and, for \(n\) sufficiently large,

\[
| \mu | (\mathrm{L} _ {n} \cup \mathrm{H}) = | \mu | (\mathrm{L} _ {n}) + | \mu | (\mathrm{H}) > \alpha .
\]

But \(L_{n} \cup H\) belongs to \(\mathfrak{K}\) by \((\mathrm{PL}_{\mathrm{II}})\) , and this contradicts the definition of \(\alpha\) .

DEFINITION 6. --- Let A be a \(\mu\) -measurable subset of X. A set \(\mathfrak{K}\) of compact subsets of A is said to be \(\mu\) -dense in A if it satisfies the conditions \((\mathrm{PL}_{\mathrm{I}})\) , \((\mathrm{PL}_{\mathrm{II}})\) , \(a)\) , \(b)\) , \(c)\) , \(d)\) of Prop. 12.

The set of all compact subsets of A is \(\mu\) -dense in A.

When A = X, we shall say simply ` \(\mu\) -dense set' in place of `set \(\mu\) -dense in X'. If X - A is locally \(\mu\) -negligible, then every set of compact subsets of A that is \(\mu\) -dense in A is also \(\mu\) -dense in X.

Remark. --- Suppose that A is the union of a sequence \((\mathrm{L}_{n})\) of compact sets and a \(\mu\) -negligible (resp. locally \(\mu\) -negligible) set, and let \(\mathfrak{K}\) be a set of compact subsets that is \(\mu\) -dense in A. Applying to each \(L_{n}\) the property c) of the statement of Prop. 12, we see that A is the union a sequence of compact sets belonging to \(\mathfrak{K}\) and a \(\mu\) -negligible (resp. locally \(\mu\) -negligible) set.

If K is a compact subset of X, it comes to the same to say that a set of compact subsets of K is \(\mu\) -dense in K or that it is \(\mu_{K}\) -dense in K; this follows from Lemmas 2 and 3 of No.~7 and condition b) of Prop. 12.

PROPOSITION 13. --- Let A be a \(\mu\) -measurable subset of X, \(\mathfrak{K}\) a set of compact subsets that is \(\mu\) -dense in A. Let \(\mathfrak{H}\) be a set of compact subsets of A satisfying (PL \(_{\text{I}}\) ) and (PL \(_{\text{II}}\) ) and such that, for every K \(\in\) \(\mathfrak{K}\) , the set of H \(\in\) \(\mathfrak{H}\) such that H \(\subset\) K is \(\mu_{\text{K}}\) -dense (or, what amounts to the same, \(\mu\) -dense) in K. Then \(\mathfrak{H}\) is \(\mu\) -dense in A.

For, let L be a compact subset of A. For every \(\varepsilon > 0\) there exists a \(K \in \mathfrak{K}\) such that \(K \subset L\) and \(|\mu|(\mathrm{L} - \mathrm{K}) \leqslant \varepsilon/2\) , and then an \(H \in \mathfrak{H}\) such that \(H \subset K\) and \(|\mu|(\mathrm{K} - \mathrm{H}) \leqslant \varepsilon/2\) ; it follows that \(|\mu|(\mathrm{L} - \mathrm{H}) \leqslant \varepsilon\) , whence the proposition.

\subsection{Locally countable partitions}

DEFINITION 7. --- A set \(\mathfrak{A}\) of subsets of a topological space \(T\) is said to be locally countable if, for every \(t \in T\) , there exists a neighborhood \(V\) of \(t\) such that the set of \(A \in \mathfrak{A}\) that intersect \(V\) is countable.

If the set A of subsets of T is locally countable then, for every compact subset K of T, the set of \(A \in A\) that intersect K is countable, since K can be covered by a finite number of open neighborhoods of points of K, each of which intersects only a countable set of subsets belonging to A.

Def. 7 shows that the union of a locally countable set of \(\mu\) -measurable (resp. locally \(\mu\) -negligible) subsets of a locally compact space is \(\mu\) -measurable (resp. locally \(\mu\) -negligible) (No.~1, Prop. 3 and No.~2, Prop. 5).

PROPOSITION 14. --- Let X be a locally compact space, \(\mu\) a measure on X, A a \(\mu\) -measurable subset of X, and \(\mathfrak{K}\) a set of compact subsets of A that is \(\mu\) -dense in A. There exists a locally countable set \(\mathfrak{H} \subset \mathfrak{K}\) , formed of pairwise disjoint sets, such that \(A - \bigcup_{K \in \mathfrak{H}} K\) is locally \(\mu\) -negligible and such that, for every \(K \in \mathfrak{H}\) , the support of \(\mu_K\) is all of K.

Consider sets \(\mathfrak{L} \subset \mathfrak{K}\) formed of pairwise disjoint sets such that, for every \(L \in \mathfrak{L}\) , \(\operatorname{Supp}(\mu_L) = L\) . The sets \(\mathfrak{L}\) form a subset \(\mathscr{H}\) of \(\mathfrak{P}(\mathfrak{K})\) that is nonempty (because it contains the element \(\varnothing\) ) and which we shall order by the relation of inclusion in \(\mathfrak{P}(\mathfrak{K})\) . It is immediate that \(\mathscr{H}\) is inductive; let \(\mathfrak{H}\) be a maximal element of \(\mathscr{H}\) (S, R, §6, No.~10). Let us first show that \(\mathfrak{H}\) is locally countable. Indeed, for every \(x \in X\) , let \(V\) be a relatively compact open neighborhood of \(x\) ; if \((K_i)_{1 \leqslant i \leqslant n}\) is a finite family of distinct sets of \(\mathfrak{H}\) that intersect \(V\) , then

\[
\sum_ {i = 1} ^ {n} | \mu | \left(\mathrm{K} _ {i} \cap \mathrm{V}\right) = | \mu | \left(\mathrm{V} \cap \left(\bigcup_ {i = 1} ^ {n} \mathrm{K} _ {i}\right)\right)
\]

because the \(K_{i}\) are pairwise disjoint, whence \(\sum_{i=1}^{n}| \mu |(\mathrm{K}_{i} \cap \mathrm{V}) \leqslant |\mu |(\mathrm{V})\) . It follows that if \(H_{V}\) is the set of \(K \in \mathfrak{H}\) that intersect V, then

\[
\sum_ {\mathrm{K} \in \mathfrak {H} _ {\mathrm{V}}} | \mu | (\mathrm{K} \cap \mathrm{V}) <   + \infty ,
\]

and since \(|\mu|(\mathrm{K} \cap \mathrm{V}) > 0\) for every \(\mathrm{K} \in \mathfrak{H}_{\mathrm{V}}\) , \(\mathfrak{H}_{\mathrm{V}}\) is necessarily countable. Next, let us prove that \(\mathrm{N} = \mathrm{A} - \bigcup_{\mathrm{K} \in \mathfrak{H}} \mathrm{K}\) is locally \(\mu\) -negligible. We saw above that \(\mathrm{N}\) is \(\mu\) -measurable. If \(\mathrm{N}\) were not locally negligible, it would contain a non-negligible compact set \(\mathrm{L}_0\) , hence (No.~8, Prop. 12) a non-negligible compact set \(\mathrm{L} \subset \mathrm{L}_0\) belonging to \(\mathfrak{K}\) . Since \(|\mu_{\mathrm{L}}|(\mathrm{L}) = |\mu|(\mathrm{L}) > 0\) (No.~7, Lemmas 2 and 3), the measure \(\mu_{\mathrm{L}}\) induced on \(\mathrm{L}\) by \(\mu\) is nonzero; its support \(\mathrm{S}\) is therefore a nonempty compact set belonging to \(\mathfrak{K}\) by \((\mathrm{PL}_{\mathrm{I}})\) , and \(\operatorname{Supp}(\mu_{\mathrm{S}}) = \mathrm{S}\) (No.~7, Lemma 2, (iii)). It follows that the set \(\mathfrak{H} \cup \{\mathrm{S}\}\) belongs to \(\mathscr{H}\) , which contradicts the definition of \(\mathfrak{H}\) ; the set \(\mathrm{N}\) is therefore locally negligible, which completes the proof.

\subsection{Measurable functions defined on a measurable subset}

PROPOSITION 15. --- Let X be a locally compact space, \(\mu\) a measure on X, A a \(\mu\) -measurable subset of X, and f a mapping of A into a topological space F. The following conditions are equivalent:

\begin{enumerate}
\def\labelenumi{\alph{enumi})}
\item
  The set \(\mathfrak{H}\) of compact subsets \(K\) of \(A\) such that the restriction of \(f\) to \(K\) is continuous, is \(\mu\) -dense in \(A\) .
\item
  There exists a set \(\mathfrak{K}\) of compact subsets of A, \(\mu\) -dense in A, such that the restriction of \(f\) to every \(\mathrm{K} \in \mathfrak{K}\) is \(\mu_{\mathrm{K}}\) -measurable.
\item
  There exist a homeomorphism \(j\) of \(\mathrm{F}\) onto a subspace of a topological space \(\mathrm{G}\) and a \(\mu\) -measurable mapping \(g\) of \(\mathrm{X}\) into \(\mathrm{G}\) , such that \(g|_{\mathrm{A}} = j \circ f\) .
\item
  Every extension of f to a mapping of X into F, constant on X - A, is \(\mu\) -measurable.
\end{enumerate}

It is clear that \(a)\) implies \(b)\) and that \(d)\) implies \(c)\) . The fact that \(c)\) implies \(a)\) follows from condition \(c)\) of Prop. 12 of No.~8. On the other hand, \(b)\) implies \(a)\) : for, Def. 1 shows that, for each \(K \in \mathfrak{K}\) , the set of subsets \(H \in \mathfrak{H}\) contained in \(K\) is \(\mu_{K}\) -dense in \(K\) (No.~8, Prop. 12, \(c\) ), and Prop. 13 of No.~8 shows that \(\mathfrak{H}\) is \(\mu\) -dense in \(A\) . It remains to see that \(a)\) implies \(d)\) . Let \(g\) be an extension of \(f\) to \(X\) , constant on \(X - A\) . For every compact subset \(L\) of \(X\) , \(L \cap A\) and \(L \cap (X - A)\) are \(\mu\) -integrable; therefore, for every \(\varepsilon > 0\) , there exist a compact subset \(P \subset L \cap A\) and a compact subset \(Q \subset L \cap (X - A)\) such that

\[
| \mu | ((L \cap A) - P) \leqslant \varepsilon / 4 \quad \text { and } \quad | \mu | ((L \cap (X - A)) - Q) \leqslant \varepsilon / 4.
\]

On the other hand, there exists a set \(H \in \mathfrak{H}\) contained in P such that \(|\mu|(P - H) \leqslant \varepsilon/2\) ; the restriction of g to the compact set \(K = H \cup Q\) is then continuous (g being constant on Q) and \(|\mu|(L - K) \leqslant \varepsilon\) , which completes the proof.

DEFINITION 8. --- Let X be a locally compact space, \(\mu\) a measure on X, and A a \(\mu\) -measurable subset of X. A mapping f of A into a topological space F is said to be \(\mu\) -measurable if it satisfies the equivalent conditions of Prop. 15.

If A is locally \(\mu\) -negligible, then every mapping of A into F is therefore \(\mu\) -measurable.

COROLLARY 1. --- Let X be a locally compact space, \(\mu\) a measure on X, A a \(\mu\) -measurable subset of X, and \(f\) a \(\mu\) -measurable mapping of A into a topological space F. Let \(\mathfrak{K}\) be a set of compact subsets of X, \(\mu\) -dense in X. Then, there exists a partition of A formed by a locally negligible set N and a locally countable family \((\mathrm{K}_{\lambda})_{\lambda \in \mathrm{L}}\) of sets \(\mathrm{K}_{\lambda} \in \mathfrak{K}\) , such that \(f|_{\mathrm{K}_{\lambda}}\) is continuous for every \(\lambda \in \mathrm{L}\) .

In view of No.~9, Prop. 14, it suffices to show that the set \(\mathfrak{H} \subset \mathfrak{K}\) of subsets \(K \in \mathfrak{K}\) such that \(K \subset A\) and \(f|K\) is continuous, is \(\mu\) -dense in \(A\) . Now, it follows at once from Prop. 1 of No.~1 and condition \(d)\) of Prop. 15 that, for every compact subset \(K_0\) of \(A\) and every \(\varepsilon > 0\) , there exists a subset \(K \subset K_0\) belonging to \(\mathfrak{K}\) such that \(|\mu|(K_0 - K) \leqslant \varepsilon\) and \(f|K\) is continuous; the conclusion therefore follows from Prop. 12 of No.~8.

COROLLARY 2. --- Let K be a compact subspace of X; in order that a mapping f of K into a topological space F be \(\mu\) -measurable, it is necessary and sufficient that it be \(\mu_{K}\) -measurable.

In view of Lemma 2 of No.~7, this follows at once from Prop. 1 of No.~1 and condition \(a\) ) of Prop. 15.

PROPOSITION 16. --- Let A be a locally countable set of μ-measurable subsets of X and let \(B = \bigcup_{A \in \mathfrak{A}} A\) . For a mapping f of B into a topological space F to be μ-measurable, it is necessary and sufficient that the restriction of f to every \(A \in \mathfrak{A}\) be μ-measurable.

We have already observed (No.~9) that B is \(\mu\) -measurable. The condition being obviously necessary, let us prove that it is sufficient. Thus, let K be a compact subset of B. By hypothesis, there exists a sequence \((\mathbf{A}_{n})\) of sets belonging to \(\mathfrak{A}\) such that the \(K \cap A_{n}\) form a covering of K. Set \(C_{0} = K \cap A_{0}\) and

\[
\mathrm{C} _ {n} = \mathrm{K} \cap \mathrm{A} _ {n} \cap \mathbf {C} \left(\bigcup_ {i <   n} \mathrm{C} _ {i}\right)
\]

for n \textgreater{} 0, so that the nonempty \(C_{n}\) form a partition of K into \(\mu\) -integrable sets. Since the restriction of f to \(C_{n}\) is \(\mu\) -measurable, there exists a partition of \(C_{n}\) formed by a \(\mu\) -negligible set \(N_{n}\) and a sequence \((\mathrm{L}_{mn})_{m \geqslant 0}\) of compact sets such that \(f|L_{mn}\) is continuous. Since \(N = \bigcup_{n} N_{n}\) is \(\mu\) -negligible, we see that condition a) of Prop. 15 is satisfied, whence the proposition.

Property d) of Prop. 15 makes it possible to immediately generalize the properties of measurable functions defined on all of X, observed in Nos. 2 to 5, to measurable functions defined on a measurable subset A of X; these generalizations are left to the reader. We only point out explicitly that the principle of localization (No.~2, Prop. 4) remains valid when it is assumed that each of the functions \(g_{x}\) is only defined in \(V_{x}\) (or almost everywhere in \(V_{x}\) ) and is measurable.

\subsection{Convergence in measure}

Let X be a locally compact space, \(\mu\) a measure on X, A a \(\mu\) -measurable subset of X, and F a uniform space; we shall denote by \(\mathcal{S}(\mathrm{A},\mu;\mathrm{F})\) , or \(\mathcal{S}_{\mathrm{F}}(\mathrm{A},\mu)\) (or simply \(\mathcal{S}_{\mathrm{F}}(\mu)\) , or even \(S_{F}\) , when A = X) the set of \(\mu\) -measurable mappings of A into F (No.~10, Def. 8). For every entourage V of the uniform structure of F, every \(\mu\) -integrable set \(B \subset A\) and every number \(\delta > 0\) , we shall denote by \(\boldsymbol{W}(\mathrm{V},\mathrm{B},\delta)\) the set of pairs \((f,g)\) of functions in \(\mathcal{S}(\mathrm{A},\mu;\mathrm{F})\) having the following property: the set M of all \(x \in B\) for which \((f(x),g(x)) \notin V\) is such that \(|\mu|^{*}(\mathrm{M}) \leqslant \delta\) . Let us show that the sets \(\boldsymbol{W}(\mathrm{V},\mathrm{B},\delta)\) form a fundamental system of entourages for a uniform structure on \(\mathcal{S}(\mathrm{A},\mu ;\mathrm{F})\) : it is clear that \(W(\mathrm{V},\mathrm{B},\delta)\) is symmetric if V is, and that if \(\mathrm{V}'\subset \mathrm{V}\) , \(\mathrm{B}'\supset \mathrm{B}\) and \(\delta^{\prime}\leqslant \delta\) , then

\[
\boldsymbol {W} \left(\mathrm{V} ^ {\prime}, \mathrm{B} ^ {\prime}, \delta^ {\prime}\right) \subset \boldsymbol {W} (\mathrm{V}, \mathrm{B}, \delta);
\]

it therefore suffices to verify the axiom \((\mathrm{U}_{\mathrm{III}}^{\prime})\) (GT, II, §1, No.~1). Now, if \(V'\) is an entourage such that \(V^{\prime} \subset V\) , then

\[
\boldsymbol {W} \left(\mathrm{V} ^ {\prime}, \mathrm{B}, \delta / 2\right) \circ \boldsymbol {W} \left(\mathrm{V} ^ {\prime}, \mathrm{B}, \delta / 2\right) \subset \boldsymbol {W} (\mathrm{V}, \mathrm{B}, \delta).
\]

Note that as K runs over a \(\mu\) -dense set \(\mathfrak{K}\) of compact subsets of A, the sets \(\boldsymbol{W}(\mathrm{V},\mathrm{K},\delta)\) also form a fundamental system of entourages for the preceding uniform structure: for, for every integrable set \(B\subset A\) , there exists a compact set \(K\in \mathfrak{K}\) contained in B such that \(|\mu|(\mathrm{B}-\mathrm{K})\leqslant\delta\) , and therefore \(\boldsymbol{W}(\mathrm{V},\mathrm{K},\delta)\subset\boldsymbol{W}(\mathrm{V},\mathrm{B},2\delta)\) .

DEFINITION 9. --- The uniform structure on \(\mathcal{S}(\mathrm{A},\mu;\mathrm{F})\) of which the \(\boldsymbol{W}(\mathrm{V},\mathrm{B},\delta)\) form a fundamental system of entourages is called the uniform structure of convergence in measure in A.

The corresponding topology is called the topology of convergence in measure in A, and a filter (or a sequence) that converges for this topology is said to be convergent in measure in A; the mention of A is often suppressed when A = X.

Suppose F is Hausdorff; then, for every \(\mu\) -integrable set \(B \subset A\) , the intersection of the entourages \(\boldsymbol{W}(V, B, \delta)\) , where V runs over a fundamental system of entourages of F and \(\delta\) runs over the set of numbers \textgreater0, is the set of pairs \((f, g)\) such that \(f(x) = g(x)\) almost everywhere in B (with respect to \(\mu\) ). For, the set M of \(x \in B\) such that \(f(x) \neq g(x)\) is \(\mu\) -integrable, because it is the inverse image, under the \(\mu\) -measurable mapping \(x \mapsto (f(x), g(x))\) , of the complement of the diagonal in \(F \times F\) , which is open (No.~5, Prop. 7); if \(|\mu|(M) = \alpha > 0\) , there exists a compact subset \(K \subset M\) such that \(|\mu|(M - K) < \alpha/2\) and such that the restrictions of f and g to K are continuous; therefore, there exists an entourage \(V_0\) of F such that \((f(x), g(x)) \notin V_0\) for all \(x \in K\) , consequently \((f, g) \notin W(V_0, B, \alpha/2)\) .

From this it follows that if F is Hausdorff, then the intersection of all the entourages of \(\mathcal{S}(\mathrm{A},\mu;\mathrm{F})\) is the set of pairs \((f,g)\) such that \(f(x)=g(x)\) locally almost everywhere in A. The Hausdorff uniform space associated with \(\mathcal{S}(\mathrm{A},\mu;\mathrm{F})\) , which we shall denote \(\mathrm{S}(\mathrm{A},\mu;\mathrm{F})\) or \(\mathrm{S}_{\mathrm{F}}(\mathrm{A},\mu)\) (or even \(\mathrm{S}_{\mathrm{F}}(\mu)\) or \(S_{F}\) when A=X), therefore consists of the equivalence classes for the relation \(\langle f(x)=g(x)\) locally almost everywhere in A \(\rangle\) in the set \(\mathcal{S}(\mathrm{A},\mu;\mathrm{F})\) .

PROPOSITION 17. --- Let \((\mathrm{A}_{\lambda})_{\lambda \in \mathrm{L}}\) be a locally countable family of \(\mu\) -measurable subsets of A, pairwise disjoint and such that \(\mathrm{A} - \bigcup_{\lambda \in \mathrm{L}} \mathrm{A}_{\lambda}\) is locally \(\mu\) -negligible. If, for every class \(\dot{f} \in \mathrm{S}(\mathrm{A}, \mu; \mathrm{F})\) and every \(\lambda \in \mathrm{L}\) , \(\dot{f}_{\lambda}\) denotes the class of the restriction to \(\mathrm{A}_{\lambda}\) of any of the functions in the class \(\dot{f}\) , then the mapping \(\psi: \dot{f} \mapsto (\dot{f}_{\lambda})_{\lambda \in \mathrm{L}}\) is an isomorphism of the uniform space \(\mathrm{S}(\mathrm{A}, \mu; \mathrm{F})\) onto the product uniform space \(\prod_{\lambda \in \mathrm{L}} \mathrm{S}(\mathrm{A}_{\lambda}, \mu; \mathrm{F})\) .

It follows from No.~10, Prop. 16 that \(\psi\) is bijective. Consider an entourage T of S(A, \(\mu\) ; F) that is the canonical image of a \(W(V, B, \delta)\) , where B is a compact subset of A; we know that the set J of \(\lambda \in L\) such that B \(\cap\) A \(_{\lambda}\) \(\neq\) \(\varnothing\) is countable (No.~9), and \(|\mu|(B) = \sum_{\lambda \in J} |\mu|(B \cap A_{\lambda})\) ; therefore, there exists a finite subset H of J such that

\[
\sum_ {\lambda \in \mathrm{J} - \mathrm{H}} | \mu | (\mathrm{B} \cap \mathrm{A} _ {\lambda}) \leqslant \frac {\delta}{2}.
\]

The image of T under \(\psi \times \psi\) is then contained in the canonical image of the product \(\prod_{\lambda \in H} W(V, B \cap A_{\lambda}, \delta)\) . On the other hand, if \(m\) is the number of elements of H, then the image of T under \(\psi \times \psi\) contains the canonical image of the entourage \(\prod_{\lambda \in H} W(V, A_{\lambda}, \delta / 2m)\) , which proves the proposition.

PROPOSITION 18. --- If F is metrizable, and A is the union of a locally \(\mu\) -negligible set and a sequence \((A_{n})\) of \(\mu\) -integrable sets, then the space \(S(A,\mu;F)\) is metrizable.

Since each \(A_{n}\) is the union of a negligible set and a sequence of compact sets, we can suppose that the \(A_{n}\) are already compact and pairwise disjoint. Prop. 17 then allows us to suppose that A is compact. If \((\mathrm{V}_{n})\) is a countable fundamental system of entourages of F, it is clear that the \(\boldsymbol{W}(\mathrm{V}_{n},\mathrm{A},1/n)\) form a fundamental system of entourages of \(\mathcal{S}(\mathrm{A},\mu;\mathrm{F})\) as n runs over N, whence the proposition.

Lemma 4. --- Let F be a metrizable uniform space, and let B ⊂ A be a countable union of μ-integrable sets. Then, for every Cauchy sequence ( \(f_{n}\) ) in \(\mathcal{S}(A,\mu;F)\) , there exists a subsequence ( \(f_{n_{k}}\) ) of ( \(f_{n}\) ) such that ( \(f_{n_{k}}(x)\) ) is a Cauchy sequence in F for almost every \(x \in B\) .

Suppose first that B is integrable, and denote by d a metric compatible with the uniform structure of F. We shall define recursively a double sequence \((f_{mn})\) of functions in \(\mathcal{S}(\mathrm{A},\mu;\mathrm{F})\) such that \(f_{0n}=f_{n}\) for every n, \((f_{mn})_{n\geqslant0}\) is a subsequence of \((f_{m-1,n})_{n\geqslant0}\) for every m\textgreater0 and, finally, such that for every m\textgreater0 the set \(M_{mn}\) of \(x\in B\) for which \(d(f_{mn}(x),f_{m,n+1}(x))>1/2^{m+n+1}\) has measure \(|\mu|(\mathrm{M}_{mn})\leqslant1/2^{m+n+1}\) ;

the possibility of such a definition follows from the fact that \((f_n)\) is a Cauchy sequence in \(\mathcal{S}(\mathrm{A},\mu ;\mathrm{F})\) . Set \(\mathbf{M}_m = \bigcup_{n\geqslant 0}\mathbf{M}_{mn}\) ; then

\[
| \mu | (\mathrm{M} _ {m}) \leqslant \sum_ {n = 0} ^ {\infty} | \mu | (\mathrm{M} _ {m n}) \leqslant 1 / 2 ^ {m}
\]

and, for every \(x \in \mathrm{B} - \mathrm{M}_m\) , we have \(d(f_{mn}(x), f_{m,n+p}(x)) \leqslant 1/2^{m+n}\) for all \(n \geqslant 0\) and all \(p > 0\) ; the sequence \((f_{mn}(x))_{n \geqslant 0}\) is therefore a Cauchy sequence in F. Now let \(\mathbf{N} = \bigcap_{m=0}^{\infty} \mathbf{M}_m\) ; N is negligible. Set \(g_n = f_{nn}\) for every \(n \geqslant 0\) ; for every \(m\) , the sequence \((g_n)_{n \geqslant m}\) is a subsequence of the sequence \((f_{mn})_{n \geqslant 0}\) ; if \(x \in \mathrm{B} - \mathrm{N}\) , there exists an index \(m\) such that \(x \notin \mathrm{M}_m\) , which proves that the sequence \((g_n(x))\) is a Cauchy sequence in F.

If now B is the union of a sequence \((\mathrm{B}_{m})\) of integrable sets, one can define recursively a double sequence \((g_{mn})\) such that \(g_{0n} = f_{n}\) , \((g_{mn})_{n \geqslant 0}\) is a subsequence of \((g_{m-1,n})_{n \geqslant 0}\) for every m \textgreater{} 0, and such that the sequence \((g_{mn}(x))_{n \geqslant 0}\) is a Cauchy sequence in \(B_{m} - P_{m}\) , where \(P_{m}\) is negligible. Set \(h_{n} = g_{nn}\) for every \(n \geqslant 0\) , so that, for every m, the sequence \((h_{n})_{n \geqslant m}\) is a subsequence of \((g_{mn})_{n \geqslant 0}\) ; the sequence \((h_{n}(x))_{n \geqslant 0}\) is then a Cauchy sequence in F for every \(x \in B - P\) , where \(P = \bigcup_{m=0}^{\infty} P_{m}\) is negligible.

PROPOSITION 19. --- If the uniform space F is metrizable and complete, then the uniform space S(A, μ; F) is complete.

There exists a locally countable family \((\mathrm{K}_{\lambda})_{\lambda \in L}\) of compact subsets of A such that the \(\mathbf{K}_{\lambda}\) are pairwise disjoint and \(\mathbf{A} - \bigcup_{\lambda}\mathbf{K}_{\lambda}\) is locally negligible (No.~9, Prop. 14). By Prop. 17, \(\mathrm{S}(A,\mu;F)\) is isomorphic to the product \(\prod_{\lambda \in L}\mathrm{S}(K_{\lambda},\mu;F)\) ; we are thus reduced to proving the proposition when A is integrable; \(\mathrm{S}(A,\mu;F)\) is then metrizable (Prop. 18) and, by Lemma 4, for every Cauchy sequence \((f_n)\) in \(\mathcal{S}(\mathrm{A},\mu;\mathrm{F})\) there exists a subsequence \((f_{n_k})\) that is convergent in \(\mathbf{A} - \mathbf{N}\) , where \(\mathbf{N}\) is negligible; the limit \(f\) of \((f_{n_k})\) (extended in any way whatsoever to all of A) is then \(\mu\) -measurable, and it follows from the extension of Egoroff's theorem mentioned in No.~10 that the sequence \((f_{n_k})\) converges in measure to \(f\) in A. This implies that \(f\) is a cluster point of the sequence \((f_n)\) in \(\mathcal{S}(\mathrm{A},\mu;\mathrm{F})\) , and since the sequence \((f_n)\) is by hypothesis a Cauchy sequence, it converges to \(f\) .

Q.E.D.

COROLLARY. --- Let F be a metrizable uniform space.

\begin{enumerate}
\def\labelenumi{(\roman{enumi})}
\item
  Every sequence \((f_n)\) of elements of \(\mathcal{S}(\mathrm{A},\mu ;\mathrm{F})\) that converges locally almost everywhere to a mapping \(f\) (necessarily \(\mu\) -measurable) of A into F, converges in measure to \(f\) in A.
\item
  Let \((f_n)\) be a sequence of elements of \(\mathcal{S}(\mathrm{A},\mu ;\mathrm{F})\) that converges in measure to a mapping \(f\) of \(\mathrm{A}\) into \(\mathrm{F}\) . For every set \(\mathbf{B}\subset \mathbf{A}\) that is a countable union of integrable sets, there exists a subsequence \((f_{n_k})\) of \((f_n)\) such that the sequence \((f_{n_k}(x))\) converges in \(\mathbf{F}\) to \(f(x)\) for almost every \(x\in \mathbf{B}\) .
\setcounter{enumi}{0}
\item
  The assertion follows at once from the extension of Egoroff's theorem mentioned in No.~10.
\item
  By Lemma 4, there exists a subsequence \((f_{n_k})\) of \((f_n)\) such that \((f_{n_k}(x))\) is a Cauchy sequence in \(\mathbf{F}\) for every \(x\in \mathrm{B} - \mathrm{N}\) , where \(\mathbf{N}\) is negligible; let \(f^{\prime}(x)\in \widehat{\mathrm{F}}\) be the limit of this sequence for \(x\in \mathrm{B} - \mathrm{N}\) . It is clear that \(f^{\prime}\) is a \(\mu\) -measurable mapping of \(\mathrm{B} - \mathrm{N}\) into \(\widehat{\mathrm{F}}\) , and the sequence \((f_n)\) converges in measure to \(f^{\prime}\) in \(\mathrm{B} - \mathrm{N}\) by (i); \(f^{\prime}\) is therefore equal to \(f\) almost everywhere in \(\mathrm{B}\) .
\end{enumerate}

PROPOSITION 20. --- Let F be a Banach space, equipped with the uniform structure defined by its norm.

\begin{enumerate}
\def\labelenumi{(\roman{enumi})}
\item
  For every \(\mu\) -measurable subset A of X, the topology of convergence in measure is compatible with the vector space structure of \(\mathcal{S}(\mathrm{A},\mu;\mathrm{F})\) .
\item
  The space \(\mathcal{K}(\mathrm{X};\mathrm{F})\) is dense in \(\mathcal{S}(\mathrm{X},\mu ;\mathrm{F})\)
\item
  For every real number \(p \geqslant 1\) , the topology induced on the space \(\mathcal{L}_{\mathrm{F}}^{p}(\mathrm{X},\mu)\) by the topology of convergence in measure is coarser than the topology of convergence in mean of order \(p\) .
\setcounter{enumi}{0}
\item
  For every \(\mu\) -integrable subset B of A and every \(\delta > 0\) , denote by \(T(B, \delta)\) the set of \(\mathbf{f} \in \mathcal{S}(A, \mu; F)\) for which the set C of \(x \in B\) such that \(|\mathbf{f}(x)| \geqslant \delta\) satisfies the relation \(|\mu|(C) \leqslant \delta\) ; it is clear that if \(V_{\delta}\) is the entourage of F formed by the pairs \((\mathbf{y}, \mathbf{z})\) such that \(|\mathbf{y} - \mathbf{z}| \leqslant \delta\) , the entourage \(W(V_{\delta}, B, \delta)\) is the set of pairs \((\mathbf{f}, \mathbf{g})\) of measurable mappings of A into F such that \(\mathbf{f} - \mathbf{g} \in T(B, \delta)\) . It is clear that the sets \(T(B, \delta)\) are symmetric, and that \(T(B, \delta) + T(B, \delta) \subset T(B, 2\delta)\) and \(T(B, |\alpha| \delta) \subset \alpha T(B, \delta)\) for every nonzero scalar \(\alpha\) such that \(|\alpha| \leqslant 1\) ; it therefore suffices to verify that the sets \(T(B, \delta)\) are absorbent (TVS, I, §1, No.~5, Prop. 4). Now, if \(\mathbf{f}\) is a \(\mu\) -measurable mapping of A into F, the numerical function \(|\mathbf{f}|\) is also \(\mu\) -measurable (No.~3, Cor. 6 of Th. 1). Let \(C_n\) be the set of \(x \in B\) such that \(|\mathbf{f}(x)| \geqslant n\) ; the \(C_n\) form a decreasing sequence of integrable sets whose intersection is empty; therefore there exists an integer \(n\) such that \(|\mu|(C_n) \leqslant \delta\) ( \(\S 4\) , No.~5, Cor. of Prop. 7); we can moreover suppose that \(n\) is taken sufficiently large that \(1/n \leqslant \delta\) ; then \(f/n^2 \in T(B, \delta)\) , which completes the proof of assertion (i).
\setcounter{enumi}{2}
\item
  The relation \(\int |\mathbf{f}|^p d|\mu| \leqslant \delta^{p+1}\) implies that if C is the set of \(x \in \mathbf{X}\) such that \(|\mathbf{f}(x)| \geqslant \delta\) , then
\end{enumerate}

\[
\delta^ {p} | \mu | ^ {*} (\mathrm{C}) \leqslant \int | \mathbf {f} | ^ {p} d | \mu | \leqslant \delta^ {p + 1},
\]

whence \(|\mu|^*(\mathrm{C}) \leqslant \delta\) , which proves (iii).

\begin{enumerate}
\def\labelenumi{(\roman{enumi})}
\setcounter{enumi}{1}
\tightlist
\item
  In view of (iii), it suffices to show for example that \(\mathcal{L}_{\mathrm{F}}^{1}\) is dense in \(\mathcal{S}_{\mathrm{F}}\) , since by definition \(\mathcal{K}(\mathrm{X};\mathrm{F})\) is dense in \(\mathcal{L}_{\mathrm{F}}^{1}\) for the topology of convergence in mean. Now, let \(\mathbf{f}\) be any element of \(\mathcal{S}_{\mathrm{F}}\) and let \(T(\mathrm{B},\delta)\) be a neighborhood of 0 in this space; we see as in (i) that there exists an integrable subset C of B such that \(|\mu|(\mathrm{C})\leqslant \delta\) and such that \(\mathbf{f}\) is bounded on \(\mathrm{B - C}\) ; denoting then by \(\mathbf{g}\) the function equal to \(\mathbf{f}\) on \(\mathrm{B - C}\) and to 0 on \(\mathrm{X - (B - C)}\) , it follows from No.~6, Th. 5 that \(\mathbf{g}\) is integrable, and obviously \(\mathbf{f} - \mathbf{g}\in T(\mathrm{B},\delta)\) .
\end{enumerate}

Remarks. --- 1) The topological vector space \(\mathcal{S}(\mathrm{X},\mu ;\mathrm{F})\) is not necessarily locally convex (Exer. 24).

\begin{enumerate}
\def\labelenumi{\arabic{enumi})}
\setcounter{enumi}{1}
\tightlist
\item
  The topology induced on the set of f such that \(\mathrm{N}_{p}(\mathbf{f}) \leqslant 1\) by the topology of convergence in measure may be strictly coarser than the topology induced on this set by the topology of convergence in mean of order p (Exer. 22). However, see Prop. 21 below.
\end{enumerate}

DEFINITION 10. --- Let X be a locally compact space, \(\mu\) a measure on X, F a Banach space, and \(p \in [1, +\infty[\) . A subset H of \(\mathcal{L}_{\mathrm{F}}^{p}(\mathrm{X}, \mu)\) is said to be equi-integrable of order \(p\) (for \(\mu\) ) if it satisfies the following conditions:

\begin{enumerate}
\def\labelenumi{(\roman{enumi})}
\tightlist
\item
  For every \(\varepsilon > 0\) there exists a \(\delta > 0\) such that, for every integrable set A of measure \(|\mu| (\mathrm{A}) \leqslant \delta\) and every \(\mathbf{f} \in \mathrm{H}\) ,
\end{enumerate}

\[
\int | \mathbf {f} | ^ {p} \varphi_ {\mathrm{A}} d | \mu | \leqslant \varepsilon .
\]

\begin{enumerate}
\def\labelenumi{(\roman{enumi})}
\setcounter{enumi}{1}
\tightlist
\item
  For every \(\varepsilon > 0\) there exists a compact subset \(K\) of \(X\) such that, for every \(\mathbf{f} \in H\) , \(\int |\mathbf{f}|^p \varphi_{X - K} d|\mu| \leqslant \varepsilon\) .
\end{enumerate}

When \(p = 1\) we say `equi-integrable' instead of `equi-integrable of order 1'.

Remark. --- Suppose \(\mu\) is bounded. For every a \textgreater{} 0, the set of measurable mappings of X into F such that \(|\mathbf{f}(x)| \leqslant a\) almost everywhere is equi-integrable of order p, and this is true for any \(p \in [1, +\infty[\) .

PROPOSITION 21. --- Let H be a subset of \(\mathcal{L}_{\mathrm{F}}^{p}(\mathrm{X},\mu)\) that is equi-integrable of order \(p\) . On H, the uniform structure of convergence in measure is equal to the uniform structure induced by that of \(\mathcal{L}_{\mathrm{F}}^{p}(\mathrm{X},\mu)\) .

Let \(\varepsilon > 0\) . There exist \(\delta\) and \(K\) with the properties (i) and (ii) of Def. 10. Let \(\mathbf{f}, \mathbf{g}\) in \(H\) be such that

\[
| \mathbf {f} (x) - \mathbf {g} (x) | \leqslant \left(\frac {\varepsilon}{| \mu | (\mathrm{K})}\right) ^ {1 / p}
\]

for \(x \in \mathbf{K}\) , except on a set \(\mathbf{M}\) of measure \(\leqslant \delta\) . Then

\[
\left(\int_ {\mathrm{X} - \mathrm{K}} | \mathbf {f} - \mathbf {g} | ^ {p} d | \mu |\right) ^ {1 / p} \leqslant \left(\int_ {\mathrm{X} - \mathrm{K}} | \mathbf {f} | ^ {p} d | \mu |\right) ^ {1 / p} + \left(\int_ {\mathrm{X} - \mathrm{K}} | \mathbf {g} | ^ {p} d | \mu |\right) ^ {1 / p}
\]

and similarly

\[
\left(\int_ {\mathrm{M}} | \mathbf {f} - \mathbf {g} | ^ {p} d | \mu |\right) ^ {1 / p} \leqslant 2 \varepsilon^ {1 / p},
\]

therefore

\[
\begin{array}{c} \int | \mathbf {f} - \mathbf {g} | ^ {p} d | \mu | = \int_ {\mathrm{X-K}} | \mathbf {f} - \mathbf {g} | ^ {p} d | \mu | + \int_ {\mathrm{M}} | \mathbf {f} - \mathbf {g} | ^ {p} d | \mu | + \int_ {\mathrm{K-M}} | \mathbf {f} - \mathbf {g} | ^ {p} d | \mu | \\ \leqslant 2 ^ {p} \varepsilon + 2 ^ {p} \varepsilon + \frac {\varepsilon}{| \mu | (\mathrm{K})} | \mu | (\mathrm{K-M}) \leqslant (2 ^ {p + 1} + 1) \varepsilon . \end{array}
\]

Thus, the uniform structure of convergence in measure on H is finer than the uniform structure induced by that of \(\mathcal{L}_{\mathrm{F}}^{p}(\mathrm{X},\mu)\) . It then suffices to apply Prop. 20.

\subsection{A property of vague convergence}

Lemma 5. --- Let X be a locally compact space, \(\mu\) a bounded positive measure on X, F a Banach space, and f a bounded function on X with values in F. The following conditions are equivalent:

\begin{enumerate}
\def\labelenumi{(\roman{enumi})}
\item
  The set of points of discontinuity of \(\mathbf{f}\) is \(\mu\) -negligible.
\item
  For every \(\varepsilon > 0\) , there exist elements \(\mathbf{a}_1, \ldots, \mathbf{a}_n\) of \(F\) , functions \(g_1, \ldots, g_n\) belonging to \(\mathcal{K}(X)\) , and a bounded continuous function \(h \geqslant 0\) on \(X\) such that \(|\mathbf{f} - g_1 \mathbf{a}_1 - \cdots - g_n \mathbf{a}_n| \leqslant h \leqslant 2 \sup |\mathbf{f}|\) on \(X\) , and \(\int h d\mu \leqslant \varepsilon\) . Denote by \(N\) the set of points of discontinuity of \(\mathbf{f}\) , and let \(M = \sup |\mathbf{f}|\) .
\setcounter{enumi}{0}
\item
  \(\Rightarrow\) (ii). Suppose that condition (i) is satisfied. Let \(\varepsilon > 0\) . The function f is \(\mu\) -integrable (No.~2, Cor. 4 of Prop. 5, and No.~6, Th. 5), therefore there exist \(a_{1}, \ldots, a_{n}\) in F and \(g_{1}, \ldots, g_{n}\) in \(\mathcal{K}(X)\) such that, on setting \(k = |f - g_{1}a_{1} - \cdots - g_{n}a_{n}|\) , we have \(\int k d\mu \leqslant \varepsilon/2\) (§3, No.~5,
Prop. 10). Multiplying \(g_1, \ldots, g_n\) by a suitable same element of \(\mathcal{K}(X)\) , we can suppose in addition that
\end{enumerate}

\[
\left| g _ {1} \mathbf {a} _ {1} + \dots + g _ {n} \mathbf {a} _ {n} \right| \leqslant \mathrm{M} = \sup | \mathbf {f} |
\]

on X, whence \(k \leqslant 2M\) . The set \(N'\) of points of discontinuity of \(k\) is contained in \(N\) , hence is negligible. For every \(x \in X\) , set \(l(x) = \limsup_{y \to x} k(y)\) .

Then \(2\mathrm{M} \geqslant l \geqslant k\) on X, and \(l = k\) on \(\mathrm{X} - \mathrm{N}'\) , that is, almost everywhere for \(\mu\) , therefore \(\int l d\mu \leqslant \varepsilon / 2\) . On the other hand, \(l\) is bounded and upper semi-continuous, hence is the lower envelope of the set of bounded continuous functions \(\geqslant l\) . Therefore there exists a bounded continuous function \(h \geqslant l\) on X such that \(h \leqslant 2\mathrm{M}\) and \(\int h d\mu \leqslant \int l d\mu + \varepsilon / 2\) (§4, No.~4, Cor. 2 of Prop. 5). Then \(\int h d\mu \leqslant \varepsilon\) and \(|\mathbf{f} - g_1\mathbf{a}_1 - \dots - g_n\mathbf{a}_n| \leqslant h\) .

\begin{enumerate}
\def\labelenumi{(\roman{enumi})}
\setcounter{enumi}{1}
\tightlist
\item
  \(\Rightarrow\) (i). Suppose that condition (ii) is satisfied. For every \(x \in X\) let \(\omega(x)\) be the oscillation of f at x (GT, IX, §2, No.~3). Let \(\varepsilon > 0\) . There exist \(a_{1}, \ldots, a_{n}, g_{1}, \ldots, g_{n}, h\) with the properties of (ii). For every \(x \in X\) , \(\omega(x)\) is the oscillation of \(f - g_{1}a_{1} - \cdots - g_{n}a_{n}\) at x, therefore \(\omega(x) \leqslant 2h(x)\) . Thus \(\int \omega d\mu \leqslant 2\varepsilon\) . Consequently, the set \(A_{\varepsilon}\) of \(x \in X\) such that \(\omega(x) \geqslant \sqrt{\varepsilon}\) satisfies \(\mu(A_{\varepsilon}) \leqslant 2\sqrt{\varepsilon}\) . This proves that \(\mu(N) \leqslant 2\sqrt{\varepsilon}\) , whence \(\mu(N) = 0\) .
\end{enumerate}

PROPOSITION 22. --- Let F be a Banach space, X a locally compact space, E the set of bounded positive measures on X, μ an element of E, and B a filter base on E. Assume that B converges vaguely to μ and that \textbar\textbar ν\textbar\textbar{} tends to \textbar\textbar μ\textbar\textbar{} with respect to B. Let f be a mapping of X into F satisfying the following conditions:

\begin{enumerate}
\def\labelenumi{(\roman{enumi})}
\item
  \(\mathbf{f}\) is bounded, and is integrable for \(\mu\) and for every measure that belongs to some element of \(\mathfrak{B}\) ;
\item
  the set of points of discontinuity of \(\mathbf{f}\) is \(\mu\) -negligible.
\end{enumerate}

Then \(\int \mathbf{f}d\nu\) tends to \(\int \mathbf{f}d\mu\) with respect to \(\mathfrak{B}\) .

Let \(\varepsilon > 0\) . There exist elements \(\mathbf{a}_1, \ldots, \mathbf{a}_n\) in \(F\) , functions \(g_1, \ldots, g_n\) in \(\mathcal{K}(X)\) , and a bounded continuous function \(h \geqslant 0\) on \(X\) , such that

\[
\left| \mathbf {f} - g _ {1} \mathbf {a} _ {1} - \dots - g _ {n} \mathbf {a} _ {n} \right| \leqslant h \leqslant 2 \sup | \mathbf {f} |
\]

on X and \(\int h d\mu \leqslant \varepsilon\) (Lemma 5). Let \(M = \sup |f|\) . There exist a compact subset K of X such that \(\mu^{*}(X - K) \leqslant \varepsilon\) (§4, No.~7, Prop. 12 and No.~6, Th. 4), a compact neighborhood \(K'\) of K in X, and a continuous mapping \(h'\) of X into {[}0,2M{]} such that \(h' = h\) on K, \(h' = 2M\) on \(X - K'\) ; replacing \(h'\) by \(\sup(h, h')\) , we can suppose in addition that \(h' \geqslant h\) . Then \(\int(h' - h) d\mu \leqslant 2M\mu^{*}(X - K) \leqslant 2M\varepsilon\) . On the other hand, \(h' = h_1 + 2M\) , where \(h_1 \in \mathcal{K}(X)\) . Taking into account §4, No.~7, Prop. 12, the number \(\int h^{\prime}d\nu = \int h_{1}d\nu +2\mathrm{M}\| \nu \| \text{ tends to } \int h_{1}d\mu +2\mathrm{M}\| \mu \| = \int h^{\prime}d\mu \text{ with respect to } \mathfrak{B}.\) There then exists an \(A\in \mathfrak{B}\) such that, for every \(\nu \in A\) ,

\[
\begin{array}{c} \Big | \int (g _ {1} \mathbf {a} _ {1} + \dots + g _ {n} \mathbf {a} _ {n}) d \nu - \int (g _ {1} \mathbf {a} _ {1} + \dots + g _ {n} \mathbf {a} _ {n}) d \mu \Big | \leqslant \varepsilon , \\ \int h d \nu \leqslant \int h ^ {\prime} d \nu \leqslant \int h ^ {\prime} d \mu + \varepsilon \leqslant \int h d \mu + 2 \mathrm{M} \varepsilon + \varepsilon \leqslant 2 (\mathrm{M} + 1) \varepsilon . \end{array}
\]

These inequalities imply

\[
\begin{array}{c} \Big | \int \mathbf {f}   d \nu - \int \mathbf {f}   d \mu \Big | \leqslant \\ \int h   d \nu + \Big | \int (g _ {1} \mathbf {a} _ {1} + \dots + g _ {n} \mathbf {a} _ {n})   d \nu - \int (g _ {1} \mathbf {a} _ {1} + \dots + g _ {n} \mathbf {a} _ {n})   d \mu \Big | + \int h   d \mu \\ \leqslant 2 (\mathrm{M} + 2) \varepsilon , \end{array}
\]

which proves the proposition.

Remark. --- The conditions (i) and (ii) of Proposition 22 are satisfied if f is continuous and bounded.

Example. --- Let us take for X the compact space U of complex numbers of absolute value 1. On setting \(\mu(f) = \int_{0}^{1} f(e^{2i\pi t}) dt\) for every \(f \in \mathcal{K}(U)\) , one defines a positive measure of mass 1 on U. On the other hand, let \(\theta\) be a real number; for every integer \(n \geqslant 0\) , let \(\nu_{n}\) be the unit mass placed at the point \(e^{2i\pi n\theta}\) of U, and let

\[
\mu_ {n} = \frac {1}{n + 1} (\nu_ {0} + \dots + \nu_ {n}),
\]

so that \(\mu_n\) is a positive measure of mass 1 on \(\mathbf{U}\) . Then, if \(\theta\) is irrational, \(\mu_n\) tends vaguely to \(\mu\) . For, since the linear combinations of the functions \(z \mapsto z^k\) ( \(k \in \mathbf{Z}\) ) are dense in \(\mathcal{K}(\mathbf{U})\) (GT, X, §4, No.~4, Prop. 8), it suffices to prove that \(\mu_n(z^k)\) tends to \(\mu(z^k)\) for \(k \in \mathbf{Z}\) . Now, for \(k = 0\) , \(\mu_n(z^k) = \mu(z^k) = 1\) ; for \(k \neq 0\) ,

\[
\mu_ {n} (z ^ {k}) = \frac {1}{n + 1} (1 + e ^ {2 i \pi k \theta} + e ^ {4 i \pi k \theta} + \dots + e ^ {2 i \pi k n \theta}).
\]

Since \(e^{2i\pi k\theta} \neq 1\) (because \(\theta\) is irrational), we infer that

\[
\left| \mu_ {n} (z ^ {k}) \right| = \left| \frac {1}{n + 1} \frac {e ^ {2 i \pi k (n + 1) \theta} - 1}{e ^ {2 i \pi k \theta} - 1} \right| \leqslant \frac {1}{n + 1} \frac {2}{\left| e ^ {2 i \pi k \theta} - 1 \right|},
\]

therefore \(\mu_n(z^k)\) tends to \(0 = \mu (z^k)\) . Under these conditions, Proposition 22 can be applied, and we see in particular that if A is a subset of U with negligible boundary with respect to \(\mu\) , then \(\mu_n(A)\) tends to \(\mu(A)\) . In other words, if \(p_n\) denotes the number of integers \(k \in [0, n]\) such that \(e^{2i\pi k\theta} \in A\) , then \(n^{-1}p_n\) tends to \(\mu(A)\) as \(n\) tends to \(+\infty\) .

\section{CONVEXITY INEQUALITIES}

\subsection{The convexity theorem}

THEOREM 1. --- Let X be a locally compact space, \(\mu\) a positive measure on X, F a real Banach space, D a closed convex set in F, and f a function on X such that \(\mathbf{f}(\mathbf{X}) \subset \mathbf{D}\) . For every non-negligible integrable numerical function \(g \geqslant 0\) such that \(\mathbf{fg}\) is integrable, the point \(\frac{\int \mathbf{fg} d\mu}{\int g d\mu}\) belongs to D.

For, let \(\mathbf{F}'\) be the dual of \(\mathbf{F}\) and let \(\langle \mathbf{z},\mathbf{a}'\rangle \leqslant \alpha (\mathbf{a}'\in \mathbf{F}',\alpha \in \mathbf{R})\) be a relation defining a closed half-space containing D. Since \(\mathbf{fg}\) is integrable, so is the numerical function \(\langle \mathbf{fg},\mathbf{a}'\rangle = \langle \mathbf{f},\mathbf{a}'\rangle g\) , and

\[
\int \left\langle \mathbf {f} g, \mathbf {a} ^ {\prime} \right\rangle d \mu = \left\langle \int \mathbf {f} g d \mu , \mathbf {a} ^ {\prime} \right\rangle
\]

(§4, No.~2, Cor. 1 of Th. 1); but, by hypothesis, \(\langle \mathbf{f}(x), \mathbf{a}' \rangle \leqslant \alpha\) for all \(x \in X\) , therefore \(\langle \mathbf{f}(x)g(x), \mathbf{a}' \rangle \leqslant \alpha g(x)\) ; on integrating, we have

\[
\left\langle \int \mathbf {f} g d \mu , \mathbf {a} ^ {\prime} \right\rangle \leqslant \alpha \int g d \mu .
\]

This proves that the point \(\frac{\int\mathbf{f}g d\mu}{\int g d\mu}\) belongs to every closed half-space containing D; but, by the Hahn--Banach theorem, D is the intersection of the closed half-spaces containing it (TVS, II, §5, No.~3, Cor. 1 of Prop. 4), whence the theorem.

COROLLARY. --- If the positive measure \(\mu\) has total mass equal to 1 and if \(\mathbf{f}\) is integrable, then \(\int \mathbf{f} d\mu\) belongs to the closed convex envelope of \(\mathbf{f}(\mathbf{X})\) in \(\mathbf{F}\) .

It suffices to take for g the constant function 1.

\subsection{Inequality of the mean}

We are going to sharpen Th. 1 for numerical measurable functions (finite or not).

DEFINITION 1. --- Let X be a locally compact space, \(\mu\) a measure on X. Given a numerical function f (finite or not), defined locally almost everywhere in X, we call maximum in measure, or \(\mu\) -maximum (resp. minimum in measure, or \(\mu\) -minimum) of the function f, and denote by \(\mathrm{M}_{\infty}(f)\) (resp. \(m_{\infty}(f)\) ), the infimum (resp. supremum) of the set of numbers \(\alpha\) such that \(f(x) \leqslant \alpha\) (resp. \(f(x) \geqslant \alpha\) ) locally almost everywhere (for \(\mu\) ).

It follows at once from the definition that \(m_{\infty}(f) = -M_{\infty}(-f)\) , thus from every property of the maximum in measure one deduces a corresponding property of the minimum in measure.

For every \(\alpha > \mathrm{M}_{\infty}(f)\) , the set of \(x \in \mathbf{X}\) such that \(f(x) > \alpha\) is locally negligible; now, the set of \(x \in \mathbf{X}\) such that \(f(x) > \mathrm{M}_{\infty}(f)\) is the union of the sets where \(f(x) > r_n\) , with \(r_n\) running over the set of rational numbers \(> \mathrm{M}_{\infty}(f)\) ; therefore \(f(x) \leqslant \mathrm{M}_{\infty}(f)\) locally almost everywhere (§5, No.~2). Similarly \(f(x) \geqslant m_{\infty}(f)\) locally almost everywhere; it follows that \(m_{\infty}(f) \leqslant \mathrm{M}_{\infty}(f)\) if the measure \(\mu\) is nonzero; moreover, the relation \(m_{\infty}(f) = \mathrm{M}_{\infty}(f)\) is equivalent to saying that \(f\) is equal to a constant locally almost everywhere. It is clear that if the measure \(\mu\) is nonzero, then

\[
\inf _ {x \in \mathrm{X}} f (x) \leqslant m _ {\infty} (f) \leqslant \mathrm{M} _ {\infty} (f) \leqslant \sup _ {x \in \mathrm{X}} f (x).
\]

If two functions \(f, g\) are equal locally almost everywhere, then \(m_{\infty}(f) = m_{\infty}(g)\) and \(\mathrm{M}_{\infty}(f) = \mathrm{M}_{\infty}(g)\) .

Finally, if \(f\) and \(g\) are two functions such that \(f + g\) is defined locally almost everywhere, then

\[
\mathrm{M} _ {\infty} (f + g) \leqslant \mathrm{M} _ {\infty} (f) + \mathrm{M} _ {\infty} (g)\tag{1}
\]

provided the second member is defined, as follows at once from Def. 1; similarly, if f and g are both \(\geqslant0\) and are such that fg is defined locally almost everywhere, then

\[
\mathrm{M} _ {\infty} (f g) \leqslant \mathrm{M} _ {\infty} (f) \mathrm{M} _ {\infty} (g)\tag{2}
\]

provided the second member is defined.

If \(M_{\infty}(f) < +\infty\) , then \(f(x) < +\infty\) locally almost everywhere, but not necessarily almost everywhere. A numerical function f is said to be bounded in measure (for the measure \(\mu\) ) if it is defined and finite almost everywhere and if, moreover, the numbers \(m_{\infty}(f)\) and \(M_{\infty}(f)\) are both finite (the latter condition amounts to saying that \(M_{\infty}(|f|)<+\infty\) ).

PROPOSITION 1 (Inequality of the mean). --- Let f be a measurable numerical function that is bounded in measure. For every integrable numerical function \(g \geqslant 0\) , the function fg (defined almost everywhere) is integrable and

\[
m _ {\infty} (f) \int g d | \mu | \leqslant \int f g d | \mu | \leqslant \mathrm{M} _ {\infty} (f) \int g d | \mu |.\tag{3}
\]

Moreover, two of the three members of the inequality (3) cannot be equal unless, in the set of \(x \in X\) such that \(g(x) \neq 0\) , \(f\) is equal to \(M_{\infty}(f)\) almost everywhere or equal to \(m_{\infty}(f)\) almost everywhere.

Indeed, \(fg\) is measurable (§5, No.~3, Cor. 5 of Th. 1); moreover, the inequality \(m_{\infty}(f)g(x) \leqslant f(x)g(x) \leqslant M_{\infty}(f)g(x)\) holds, not only locally almost everywhere, but even almost everywhere, because the set of points \(x \in X\) where \(g(x) \neq 0\) is a countable union of integrable sets (§5, No.~6, Lemma 1). It follows that \(fg\) is integrable (§5, No.~6, Th. 5) and the inequality (3) holds. On the other hand, the function \(M_{\infty}(f)g - fg\) is almost everywhere defined and equal to \((M_{\infty}(f) - f)g\) ; it is therefore \(\geqslant 0\) almost everywhere in \(X\) ; since the relation \(M_{\infty}(f) \int g d|\mu| = \int fg d|\mu|\) is equivalent to \(\int(M_{\infty}(f) - f)gd|\mu| = 0\) , it can hold only if the function \((M_{\infty}(f) - f)g\) is negligible, which completes the proof.

Setting aside the trivial case that \(\int g d|\mu| = 0\) , the inequality (3) may be deduced from Th. 1 of No.~1 applied to the interval \(D = [m_{\infty}(f), M_{\infty}(f)]\) . One can bring to Th. 1 of No.~1 the complements analogous to those of Prop. 1, that specify the case in which the point \((\int f g d\mu)/(\int g d\mu)\) belongs to the boundary of D (Exer. 2).

\subsection{\texorpdfstring{The spaces \(\mathcal{L}_{\mathrm{F}}^{\infty}\)}{The spaces L sub F to infinity}}

DEFINITION 2. --- For every mapping f of X into a Banach space F, one sets \(\mathrm{N}_{\infty}(\mathbf{f}) = \mathrm{M}_{\infty}(|\mathbf{f}|)\) ; f is said to be bounded in measure (for the measure \(\mu\) ) if \(\mathrm{N}_{\infty}(\mathbf{f})\) is finite. The set of mappings of X into F that are measurable and bounded in measure is denoted \(\mathcal{L}_{\mathrm{F}}^{\infty}(\mathrm{X}, \mu)\) (or \(\mathcal{L}_{\mathrm{F}}^{\infty}(\mu)\) , or simply \(\mathcal{L}_{F}^{\infty}\) ).

A function \(\mathbf{f}\) in \(\mathcal{L}_{\mathrm{F}}^{\infty}\) may thus be characterized by the fact that there exists a bounded measurable function equal locally almost everywhere to \(\mathbf{f}\) .

It follows immediately from (1) that

\[
\mathrm{N} _ {\infty} (\mathbf {f} + \mathbf {g}) \leqslant \mathrm{N} _ {\infty} (\mathbf {f}) + \mathrm{N} _ {\infty} (\mathbf {g});
\]

on the other hand, \(\mathrm{N}_{\infty}(\alpha\mathbf{f}) = |\alpha| \mathrm{N}_{\infty}(\mathbf{f})\) for every scalar \(\alpha\) . The set \(\mathcal{L}_{F}^{\infty}\) is therefore a linear subspace of the space of all mappings of X into F, and \(\mathrm{N}_{\infty}(\mathbf{f})\) is a semi-norm on this vector space. Let \((\mathbf{f}_{n})\) be a sequence of functions in \(\mathcal{L}_{F}^{\infty}\) that converges to \(f \in \mathcal{L}_{F}^{\infty}\) for the topology defined by the semi-norm \(\mathrm{N}_{\infty}(\mathbf{f})\) ; for every integer m, there exist a locally negligible set \(H_{m}\) and an integer \(n_{0}\) such that \(|\mathbf{f}(x) - \mathbf{f}_{n}(x)| \leqslant 1/m\) for every integer \(n \geqslant n_{0}\) and every \(x \notin H_{m}\) (every countable union of locally negligible sets being locally negligible); the union H of the \(H_{m}\) is locally negligible, and one sees that \(\mathbf{f}_{n}(x)\) tends uniformly to \(\mathbf{f}(x)\) on the complement of the locally negligible set H; the converse is immediate.

It is clear that every function equal locally almost everywhere to a function in \(\mathcal{L}_{F}^{\infty}\) belongs to \(\mathcal{L}_{F}^{\infty}\) . In particular, the locally negligible functions defined on X with values in F form a linear subspace \(\mathcal{N}_{F}^{\infty}\) of \(\mathcal{L}_{F}^{\infty}\) , characterized by the relation \(\mathrm{N}_{\infty}(\mathbf{f}) = 0\) (the closure of 0 for the topology defined by \(\mathrm{N}_{\infty}(\mathbf{f})\) ). The Hausdorff space associated with \(\mathcal{L}_{F}^{\infty}\) , that is, the quotient space \(\mathcal{L}_{\mathrm{F}}^{\infty}/\mathcal{N}_{\mathrm{F}}^{\infty}\) , is denoted \(\mathrm{L}_{\mathrm{F}}^{\infty}(\mathrm{X}, \mu)\) (or \(\mathrm{L}_{\mathrm{F}}^{\infty}(\mu)\) or \(\mathrm{L}_{\mathrm{F}}^{\infty}\) ); its topology is defined by the norm deduced from \(N_{\infty}\) by passage to the quotient; the norm of a class \(\dot{f} \in L_{F}^{\infty}\) is denoted \(\mathrm{N}_{\infty}(\dot{f})\) , or also \(\|f\|_{\infty}\) . When F = R (resp. C), we write \(\mathcal{L}^{\infty}\) and \(L^{\infty}\) in place of \(\mathcal{L}_{R}^{\infty}\) and \(L_{R}^{\infty}\) (resp. \(\mathcal{L}_{C}^{\infty}\) and \(L_{C}^{\infty}\) ) if this can cause no confusion.

PROPOSITION 2. --- The space \(\mathcal{L}_{\mathrm{F}}^{\infty}\) is complete; the space \(\mathrm{L}_{\mathrm{F}}^{\infty}\) is a Banach space.

For, let \((\mathbf{f}_{n})\) be a Cauchy sequence in \(\mathcal{L}_{F}^{\infty}\) ; for every integer n, there exists an integer \(k_{n}\) such that \(\mathrm{N}_{\infty}(\mathbf{f}_{r}-\mathbf{f}_{s})\leqslant1/n\) for \(r\geqslant k_{n}\) and \(s\geqslant k_{n}\) ; thus, there exists a locally negligible set \(A_{rs}\) such that \(|\mathbf{f}_{r}(x)-\mathbf{f}_{s}(x)|\leqslant1/n\) for all \(x\notin A_{rs}\) . If \(A_{n}\) is the union of the sets \(A_{rs}\) (for \(r\geqslant k_{n}\) and \(s\geqslant k_{n}\) ), then \(A_{n}\) is locally negligible and, for every \(x\notin A_{n}\) , \(|\mathbf{f}_{r}(x)-\mathbf{f}_{s}(x)|\leqslant1/n\) for all indices \(r\geqslant k_{n}\) , \(s\geqslant k_{n}\) . Let A be the locally negligible set formed by the union of the \(A_{n}\) , and set \(\mathbf{g}_{n}(x)=\mathbf{f}_{n}(x)\) for \(x\notin A\) , \(\mathbf{g}_{n}(x)=0\) for \(x\in A\) ; then \(g_{n}\) belongs to \(\mathcal{L}_{F}^{\infty}\) and, by the definition of A, the sequence \((\mathbf{g}_{n})\) converges uniformly on X to a function g. It follows that the function g is measurable (§5, No.~4, Th. 2); moreover, g is bounded on the set of \(x\in X\) where \(|\mathbf{g}_{k_{1}}(x)|\leqslant\mathrm{N}_{\infty}(\mathbf{g}_{k_{1}})\) and, since the complement of this set is locally negligible, g belongs to \(\mathcal{L}_{F}^{\infty}\) . It is clear that in \(\mathcal{L}_{F}^{\infty}\) , the sequence \((\mathbf{g}_{n})\) has limit g, and the same is therefore true of the sequence \((\mathbf{f}_{n})\) , since \(\mathrm{N}_{\infty}(\mathbf{f}_{n}-\mathbf{g}_{n})=0\) for all n.~The second part of the proposition may be deduced immediately from this.

Remarks. --- 1) Every bounded continuous function \(\mathbf{f}\) on \(X\) with values in \(F\) belongs to \(\mathcal{L}_{\mathrm{F}}^{\infty}\) , and

\[
\mathrm{N} _ {\infty} (\mathbf {f}) \leqslant \| \mathbf {f} \| = \sup _ {x \in \mathrm{X}} | \mathbf {f} (x) |.
\]

In order that \(\mathrm{N}_{\infty}(\mathbf{f})=\|\mathbf{f}\|\) for every bounded continuous function f, it is necessary and sufficient that the support of the measure \(\mu\) be equal to X. For, if there exists a continuous function f with negligible compact support and not identically zero, then \(\mathrm{N}_{\infty}(\mathbf{f})=0\) and \(\|f\|>0\) . Conversely, if the support of \(\mu\) is equal to X then, for every bounded continuous function f and every number \(\alpha<\|f\|\) , the set of \(x\in X\) such that \(|\mathbf{f}(x)|>\alpha\) is open and nonempty, hence has outer measure \textgreater0, which shows that \(\mathrm{N}_{\infty}(\mathbf{f})=\|\mathbf{f}\|\) .

When the support of \(\mu\) is equal to \(X\) , we may therefore identify the normed space \(\mathcal{C}^b (X;F)\) , of bounded continuous functions on \(X\) with values in \(F\) , with a subspace of the space \(\mathcal{L}_{F}^{\infty}\) . Since \(\mathcal{L}_F^\infty\) is not in general Hausdorff, the subspace \(\mathcal{C}^b (X;F)\) is not in general closed in \(\mathcal{L}_F^\infty\) , but its canonical image in \(L_F^\infty\) is a closed subspace of \(L_F^\infty\) (which can moreover be identified with \(\mathcal{C}^b (X;F)\) in the case contemplated). In general, \(\mathcal{C}^b (X;F)\) is distinct from \(L_F^\infty\) , that is, for an arbitrary bounded measurable function \(\mathbf{f}\) , there does not in general exist a continuous function \(\mathbf{g}\) equal to \(\mathbf{f}\) locally almost everywhere ( \(\S 5\) , Exer. 12). This implies that the space \(\mathcal{K}(X;F)\) of mappings of \(X\) into \(F\) , continuous with compact support, is in general not dense in \(L_F^\infty\) , whereas it is dense in each of the spaces \(L sub F to the p\) for \(1\leqslant p < +\infty\) ( \(\S 3\) , No.~4. Def. 2).

\pandocbounded{\includegraphics[keepaspectratio]{images/part002_fffac8d3.jpg}}

\begin{enumerate}
\def\labelenumi{\arabic{enumi})}
\setcounter{enumi}{1}
\tightlist
\item
  It is immediate that the topology defined by the semi-norm \(N_{\infty}\) is finer than the topology induced on \(L_{F}^{\infty}\) by the topology of convergence in measure (§5, No.~11).
\end{enumerate}

\subsection{Hölder's inequality}

In this No., p and q will denote two real numbers such that \(1 \leqslant p \leqslant +\infty\) , \(1 \leqslant q \leqslant +\infty\) , bound by the relation \(1/p + 1/q = 1\) ; thus \(q = p/(p - 1)\) if \(1 < p < +\infty\) , \(q = +\infty\) if p = 1, and q = 1 if \(p = +\infty\) ; p and q will be called conjugate exponents. Note that the relation \(1 \leqslant p \leqslant 2\) is equivalent to \(2 \leqslant q \leqslant +\infty\) ; p = q only when p and q are equal to 2.

THEOREM 2 (Hölder's inequality). --- Let \(f\) and \(g\) be two numerical functions that are finite almost everywhere and are such that \(f\) is equal almost everywhere to a function in \(\mathcal{L}^p\) and \(g\) to a function in \(\mathcal{L}^q\) . Then, the function \(fg\) (defined almost everywhere) is integrable, and

\[
\mathrm{N} _ {1} (f g) \leqslant \mathrm{N} _ {p} (f) \mathrm{N} _ {q} (g).\tag{4}
\]

Let \(f_{1}\) (resp. \(g_{1}\) ) be a function in \(\mathcal{L}^{p}\) (resp. \(\mathcal{L}^{q}\) ) to which \(f\) (resp. \(g\) ) is equal almost everywhere; \(fg\) is equal almost everywhere to the function \(f_{1}g_{1}\) , which is everywhere defined and finite, and which is measurable, being the product of two measurable functions (§5, No.~3, Cor. 5 of Th. 1). If \(1 < p < +\infty\) , Hölder's inequality for the upper integral (Ch. I, No.~3, Prop. 4) yields the inequality (4), and the relation \(N_{1}(fg) < +\infty\) then shows that \(fg\) is integrable (§5, No.~6, Th. 5). If \(p = 1\) , \(q = +\infty\) , the inequality (4) and the fact that \(fg\) is integrable are immediate consequences of the inequality of the mean (No.~2, Prop. 1); thus the theorem is proved in all cases.

COROLLARY 1. --- Let F, G, H be three Banach spaces, and let \((\mathbf{u}, \mathbf{v}) \mapsto \Phi(\mathbf{u}, \mathbf{v})\) be a continuous bilinear mapping of \(F \times G\) into H such that \(|\Phi(\mathbf{u}, \mathbf{v})| \leqslant |\mathbf{u}| \cdot |\mathbf{v}|\) . If \(f \in L_{F}^{p}\) and \(g \in L_{G}^{q}\) , then the function \(\Phi(\mathbf{f}, \mathbf{g})\) is integrable and

\[
\left| \int \Phi (\mathbf {f}, \mathbf {g}) d \mu \right| \leqslant \int | \Phi (\mathbf {f}, \mathbf {g}) | d | \mu | \leqslant \mathrm{N} _ {p} (\mathbf {f}) \mathrm{N} _ {q} (\mathbf {g}).\tag{5}
\]

For, \(\Phi(\mathbf{f},\mathbf{g})\) is measurable (§5, No.~3, Cor. 5 of Th. 1); since \(|\Phi(\mathbf{f},\mathbf{g})| \leqslant |\mathbf{f}| \cdot |\mathbf{g}|\) , the corollary follows from Th. 2 and the integrability criterion of §5, No.~6, Th. 5.

Two special cases of Cor. 1 are important in applications:

COROLLARY 2. --- Let F be a real (resp. complex) Banach space, F' its strong dual (TVS, III, §3, No.~1), and let \((\mathbf{z}, \mathbf{z}') \mapsto \langle \mathbf{z}, \mathbf{z}' \rangle\) be the canonical bilinear form on \(F \times F'\) . If \(f \in L_{F}^{p}\) and \(g \in L_{F'}^{q}\) , then the real (resp. complex) function \(\langle f, g \rangle\) is integrable and

\[
\left| \int \langle \mathbf {f}, \mathbf {g} \rangle d \mu \right| \leqslant \int | \langle \mathbf {f}, \mathbf {g} \rangle | d | \mu | \leqslant N _ {p} (\mathbf {f}) N _ {q} (\mathbf {g}).\tag{6}
\]

For, \(|\langle \mathbf{z},\mathbf{z}^{\prime}\rangle |\leqslant |\mathbf{z}|\cdot |\mathbf{z}^{\prime}|\)

When F is a real or complex Hilbert space, one knows that it can be canonically identified with its dual \(F'\) (TVS, V, §1, No.~7). Since the space \(L_{F}^{2}\) is complete, we have the following result:

COROLLARY 3. --- Let \(\mu\) be a positive measure on X, F a real (resp. complex) Hilbert space. On the space \(L_{F}^{2}\) , the symmetric (resp. Hermitian) form

\[
(\widetilde {\mathbf {f}}, \widetilde {\mathbf {g}}) \mapsto \int \langle \mathbf {f}, \mathbf {g} \rangle d \mu
\]

defines a Hilbert space structure, for which the norm is equal to \(\| \widetilde{\mathbf{f}}\| _2\)

COROLLARY 4. --- Let F be a Banach space, f a function in \(\mathcal{L}_{\mathrm{F}}^{p}\) , and g a numerical function belonging to \(\mathcal{L}^q\) ; then, the function fg is integrable and

\[
\left| \int \mathbf {f} g d \mu \right| \leqslant \int | \mathbf {f} g | d | \mu | \leqslant \mathrm{N} _ {p} (\mathbf {f}) \mathrm{N} _ {q} (g).\tag{7}
\]

COROLLARY 5. --- Let \(f_1, f_2, \ldots, f_n\) be \(n\) positive integrable functions, and let \(\alpha_1, \alpha_2, \ldots, \alpha_n\) be \(n\) numbers \(> 0\) such that \(\sum_{i=1}^{n} \alpha_i = 1\) ; under these conditions, the function \(f_1^{\alpha_1} f_2^{\alpha_2} \cdots f_n^{\alpha_n}\) is integrable and

\[
\begin{array}{l} \int f _ {1} ^ {\alpha_ {1}} f _ {2} ^ {\alpha_ {2}} \dots f _ {n} ^ {\alpha_ {n}} d | \mu | \leqslant \\ \left(\int f _ {1} d | \mu |\right) ^ {\alpha_ {1}} \left(\int f _ {2} d | \mu |\right) ^ {\alpha_ {2}} \dots \left(\int f _ {n} d | \mu |\right) ^ {\alpha_ {n}}. \end{array}\tag{8}
\]

For, the product \(f_{1}^{\alpha_{1}}f_{2}^{\alpha_{2}}\cdots f_{n}^{\alpha_{n}}\) is measurable, being the product of measurable functions (§5, No.~3, Th. 1 and its Cor. 5); since the inequality (8) is true for upper integrals (Ch. I, No.~2, Cor. of Prop. 2), the function \(f_{1}^{\alpha_{1}}f_{2}^{\alpha_{2}}\cdots f_{n}^{\alpha_{n}}\) is integrable (§5, No.~6, Th. 5), whence the corollary.

Cor. 2 of Th. 2 is sharpened by the following proposition:

PROPOSITION 3. --- Let \(\mu\) be a positive measure on X, F a real or complex Banach space, F' its strong dual, and \((\mathbf{z},\mathbf{z}')\mapsto\langle\mathbf{z},\mathbf{z}'\rangle\) the canonical bilinear form on F × F'.

\(1^{\circ}\) For every function \(\mathbf{f} \in \mathcal{L}_{\mathrm{F}}^{p}\) ( \(1 \leqslant p \leqslant +\infty\) ),

\[
\mathrm{N} _ {p} (\mathbf {f}) = \sup \left| \int \langle \mathbf {f}, \mathbf {g} \rangle d \mu \right|\tag{9}
\]

as \(\mathbf{g}\) runs over the set of functions in \(\mathcal{L}_{\mathrm{F}'}^q\) , such that \(\mathrm{N}_q(\mathbf{g}) \leqslant 1\) .

\(2^{\circ}\) For every function \(\mathbf{g} \in \mathcal{L}_{\mathrm{F}}^{q}\) ( \(1 \leqslant q \leqslant +\infty\) ),

\[
\mathrm{N} _ {q} (\mathbf {g}) = \sup \left| \int \langle \mathbf {f}, \mathbf {g} \rangle d \mu \right|\tag{10}
\]

as \(\mathbf{f}\) runs over the set of functions in \(\mathcal{L}_{\mathrm{F}}^{p}\) such that \(\mathrm{N}_p(\mathbf{f})\leqslant 1\)

Let us first prove the relation (9); we distinguish between two cases.

\begin{enumerate}
\def\labelenumi{(\roman{enumi})}
\tightlist
\item
  \(1 \leqslant p < +\infty\) . The relation (9) being trivial when \(\mathrm{N}_p(\mathbf{f}) = 0\) (because \(\mathbf{f}\) and \(\langle \mathbf{f}, \mathbf{g} \rangle\) are then negligible), we can always suppose, on multiplying \(\mathbf{f}\) by a scalar, that \(\mathrm{N}_p(\mathbf{f}) = 1\) . Suppose first that \(\mathbf{f}\) is an integrable step function, \(\mathbf{f} = \sum_{k=1}^{n} \dot{\mathbf{a}}_k \varphi_{\mathrm{A}_k}\) , where the \(\mathrm{A}_k\) are pairwise disjoint (§4, No.~9, Lemma). Thus \(\sum_{k=1}^{n} |\mathbf{a}_k|^p \mu(\mathrm{A}_k) = 1\) by hypothesis. For every \(\varepsilon > 0\) , there exists (for every index \(k\) ) a vector \(\mathbf{a}_k' \in \mathrm{F}'\) such that \(|\mathbf{a}_k'|^q = |\mathbf{a}_k|^p\) if \(p > 1\) (resp. \(|\mathbf{a}_k'| = 1\) if \(p = 1\) ) and \(\langle \mathbf{a}_k, \mathbf{a}_k' \rangle \geqslant (1 - \varepsilon) |\mathbf{a}_k| \cdot |\mathbf{a}_k'|\) (TVS, IV, §1, No.~3, Prop. 8). Setting \(\mathbf{g} = \sum_{k=1}^{n} \mathbf{a}_{k}^{\prime} \varphi_{\mathrm{A}_{k}}\) , we have \(\sum_{k=1}^{n} |\mathbf{a}_{k}^{\prime}|^{q} \mu(\mathrm{A}_{k}) = 1\) if \(p > 1\) (resp. \(\sup_{1 \leqslant k \leqslant n} |\mathbf{a}_{k}^{\prime}| = 1\) if \(p = 1\) ), thus \(\mathrm{N}_{q}(\mathbf{g}) = 1\) ; on the other hand,
\end{enumerate}

\[
\int \langle \mathbf {f}, \mathbf {g} \rangle d \mu = \sum_ {k = 1} ^ {n} \left\langle \mathbf {a} _ {k}, \mathbf {a} _ {k} ^ {\prime} \right\rangle \mu \left(\mathrm{A} _ {k}\right) \geqslant (1 - \varepsilon) \sum_ {k = 1} ^ {n} | \mathbf {a} _ {k} | \cdot | \mathbf {a} _ {k} ^ {\prime} | \mu (\mathrm{A} _ {k})
\]

and, since \(|\mathbf{a}_k'| = |\mathbf{a}_k|^{p / q} = |\mathbf{a}_k|^{p - 1}\) if \(p > 1\) (resp. \(|\mathbf{a}_k'| = 1\) if \(p = 1\) ), we have

\[
\int \langle \mathbf {f}, \mathbf {g} \rangle d \mu \geqslant (1 - \varepsilon) \sum_ {k = 1} ^ {n} | \mathbf {a} _ {k} | ^ {p} \mu (\mathrm{A} _ {k}) = (1 - \varepsilon) \mathrm{N} _ {p} (\mathbf {f}) = 1 - \varepsilon ,
\]

which proves the relation (9) in this case.

Let us pass to the case that \(\mathbf{f}\) is any element of \(\mathcal{L}_{\mathrm{F}}^{p}\) such that \(\mathrm{N}_p(\mathbf{f}) = 1\) . For every \(\varepsilon > 0\) , there exists a step function \(\mathbf{f}_1 \in \mathcal{L}_{\mathrm{F}}^p\) such that \(\mathrm{N}_p(\mathbf{f} - \mathbf{f}_1) \leqslant \varepsilon\) (§4, No.~10, Cor. 1 of Prop. 19). By what we have just seen, there exists a function \(\mathbf{g} \in \mathcal{L}_{\mathrm{F}}^q\) , such that \(\mathrm{N}_q(\mathbf{g}) = 1\) and

\[
\int \langle \mathbf {f} _ {1}, \mathbf {g} \rangle d \mu \geqslant \mathrm{N} _ {p} (\mathbf {f} _ {1}) - \varepsilon \geqslant 1 - 2 \varepsilon .
\]

Now,

\[
\int \left\langle \mathbf {f}, \mathbf {g} \right\rangle d \mu = \int \left\langle \mathbf {f} _ {1}, \mathbf {g} \right\rangle d \mu + \int \left\langle \mathbf {f} - \mathbf {f} _ {1}, \mathbf {g} \right\rangle d \mu
\]

and, by (6),

\[
\left| \int \langle \mathbf {f} - \mathbf {f} _ {1}, \mathbf {g} \rangle d \mu \right| \leqslant \mathrm{N} _ {p} (\mathbf {f} - \mathbf {f} _ {1}) \mathrm{N} _ {q} (\mathbf {g}),
\]

whence

\[
\left| \int \langle \mathbf {f}, \mathbf {g} \rangle d \mu \right| \geqslant 1 - 3 \varepsilon ,
\]

which proves (9).

\begin{enumerate}
\def\labelenumi{(\roman{enumi})}
\setcounter{enumi}{1}
\tightlist
\item
  \(p = +\infty\) . We may again restrict ourselves to the case that \(\mathrm{N}_{\infty}(\mathbf{f}) > 0\) . Let \(\alpha\) be any number such that \(0 < \alpha < \mathrm{N}_{\infty}(\mathbf{f})\) ; by hypothesis, the set of \(x \in X\) such that \(|\mathbf{f}(x)| > \alpha\) is measurable and is not locally negligible, therefore it contains a compact set K of measure \textgreater{} 0. Since f is measurable, there exists a compact set \(K_{1} \subset K\) of measure \textgreater{} 0, such that the restriction of f to \(K_{1}\) is continuous. It follows that, for every \(\varepsilon > 0\) , there exists a partition of \(K_{1}\) into a finite number of integrable sets, in each of which the oscillation of f is \(\leqslant \varepsilon\) ; at least one of these sets A has measure \(>0\) . Let \(\mathbf{a}\) be one of the values of \(\mathbf{f}\) in \(\mathbf{A}\) ; then \(|\mathbf{a}| > \alpha\) and \(|\mathbf{f}(x) - \mathbf{a}| \leqslant \varepsilon\) for all \(x \in \mathbf{A}\) . There exists a vector \(\mathbf{a}' \in \mathbf{F}'\) such that \(|\mathbf{a}'| = 1\) and \(|\langle \mathbf{a}, \mathbf{a}' \rangle| \geqslant |\mathbf{a}| - \varepsilon\) ; the function \(\mathbf{g} = \varphi_{\mathrm{A}} \cdot \mathbf{a}' / \mu(\mathbf{A})\) is integrable and \(\mathrm{N}_1(\mathbf{g}) = 1\) ; on the other hand,
\end{enumerate}

\[
\int \langle \mathbf {f}, \mathbf {g} \rangle d \mu = \frac {1}{\mu (\mathrm{A})} \int \langle \mathbf {f}, \mathbf {a} ^ {\prime} \rangle \varphi_ {\mathrm{A}} d \mu .
\]

Now, one can write

\[
\int \langle {\bf f}, {\bf a} ^ {\prime} \rangle \varphi_ {\mathrm{A}} d \mu = \langle {\bf a}, {\bf a} ^ {\prime} \rangle \mu (\mathrm{A}) + \int \langle {\bf f} - {\bf a}, {\bf a} ^ {\prime} \rangle \varphi_ {\mathrm{A}} d \mu ,
\]

and since

\[
\left| \langle \mathbf {f} - \mathbf {a}, \mathbf {a} ^ {\prime} \rangle \varphi_ {\mathrm{A}} \right| \leqslant \varepsilon \varphi_ {\mathrm{A}},
\]

we see that

\[
\left| \int \langle \mathbf {f}, \mathbf {g} \rangle d \mu \right| \geqslant | \langle \mathbf {a}, \mathbf {a} ^ {\prime} \rangle | - \varepsilon \geqslant | \mathbf {a} | - 2 \varepsilon > \alpha - 2 \varepsilon ;
\]

since \(\varepsilon\) is arbitrary and \(\alpha\) is any number \(<\mathrm{N}_{\infty}(\mathbf{f})\) , the relation (9) is also verified in this case.

One argues in exactly the same manner to prove the relation (10), on considering separately the case \(1 \leqslant q < +\infty\) and the case \(q = +\infty\) , and using the fact that for every \(\mathbf{z}' \in \mathrm{F}'\) , \(|\mathbf{z}'| = \sup_{|\mathbf{z}| \leqslant 1} |\langle \mathbf{z}, \mathbf{z}' \rangle|\) by the definition of the norm in \(\mathrm{F}'\) .

Remarks. --- 1) Let \(\mathcal{E}\) be a dense linear subspace of \(\mathcal{L}_{\mathrm{F}'}^q\) ; then the formula (9) holds when \(\mathbf{g}\) runs over the intersection of \(\mathcal{E}\) with the set B of functions in \(\mathcal{L}_{\mathrm{F}'}^q\) such that \(\mathrm{N}_q(\mathbf{g}) \leqslant 1\) . For, it suffices to observe that the interior \(\overset{\circ}{\mathrm{B}}\) of B is dense in B and that \(\overset{\circ}{\mathrm{B}} \cap \mathcal{E}\) is dense in \(\overset{\circ}{\mathrm{B}}\) , since \(\overset{\circ}{\mathrm{B}}\) is open. This remark applies in particular to the set \(\mathcal{E} = \mathcal{K}(\mathrm{X};\mathrm{F}')\) of continuous functions with compact support (with values in \(\mathrm{F}'\) ) when \(1 \leqslant q < +\infty\) , that is, \(1 < p \leqslant +\infty\) . But in this case, the formula (9) is true as \(\mathbf{g}\) runs over \(\mathrm{B} \cap \mathcal{K}(\mathrm{X};\mathrm{F}')\) , even for \(p = 1\) . For, we may, as above, restrict ourselves to the case that \(\mathbf{f}\) is a step function. We saw then that if \(\mathrm{N}_1(\mathbf{f}) = 1\) , then for every \(\varepsilon > 0\) there exists a step function \(\mathbf{g} \in \mathcal{L}_{\mathrm{F}'}^{\infty}\) such that \(|\mathbf{g}(x)| \leqslant 1\) for all \(x \in \mathrm{X}\) and \(|\int \langle \mathbf{f},\mathbf{g}\rangle d\mu| \geqslant 1 - \varepsilon\) . There exists a finite number of pairwise disjoint compact sets \(\mathrm{K}_i\) such that \(\mathbf{g}\) has a constant value \(\mathbf{a}_i'\) on each \(\mathrm{K}_i\) and such that, if K is the union of the \(\mathrm{K}_i\) , then \(\int |\mathbf{f}| \varphi_{\mathbf{C} K} d\mu \leqslant \varepsilon\) . Let \(\mathrm{U}_i\) be a neighborhood of \(\mathrm{K}_i\) such that the sets \(\mathrm{U}_i\) are pairwise disjoint, and let \(h_i\) be a continuous mapping of X into {[}0,1{]} with support contained in \(\mathrm{U}_i\) and equal to 1 on \(\mathrm{K}_i\) . Setting \(\mathbf{h} = \sum \mathbf{a}_i' h_i\) , we have \(\mathbf{h}(x) = \mathbf{g}(x)\) on K and \(|\mathbf{h}(x)| \leqslant 1\) on X, therefore

\[
\int | \langle \mathbf {f}, \mathbf {h} \rangle | \varphi_ {\mathbf {C K}} d \mu \leqslant \varepsilon
\]

and consequently \(|\int \langle \mathbf{f},\mathbf{h}\rangle d\mu |\geqslant 1 - 3\varepsilon\) , which proves our assertion. Analogous remarks can be made for the formula (10).

\begin{enumerate}
\def\labelenumi{\arabic{enumi})}
\setcounter{enumi}{1}
\tightlist
\item
  Let \(\mu\) be a positive measure on \(X\) , \(f\) a measurable function \(\geqslant 0\) (finite or not) whose support is contained in a countable union of compact sets \(K_{n}\) . Then, for every \(p\) such that \(1 \leqslant p \leqslant +\infty\) ,
\end{enumerate}

\[
\mathrm{N} _ {p} (f) = \sup \int^ {*} | f g | d \mu ,\tag{11}
\]

as \(g\) runs over the set of functions in \(\mathcal{K}(\mathrm{X};\mathbf{R})\) such that \(\mathrm{N}_q(g)\leqslant 1\) . For, the formula (11) is a special case of (9) when \(\mathrm{N}_p(f) < +\infty\) , since \(f\) is then equivalent to a function in \(\mathcal{L}^p\) (§5, No.~6, Th. 5). If \(\mathrm{N}_p(f) = +\infty\) , for every integer \(n > 0\) set \(f_{n} = \inf (n,f\varphi_{\mathrm{K}_{n}})\) . Then

\[
\mathrm{N} _ {p} (f _ {n}) = \sup \int^ {*} | f _ {n} g | d \mu \leqslant \sup \int^ {*} | f g | d \mu ,
\]

whence, on passing to the limit (assuming, as we may, that the sequence \((\mathbf{K}_n)\) is increasing), we have \(\sup \int^{*}|fg| d\mu = +\infty\) (§1, No.~3, Th. 3).

COROLLARY. --- Let \(\mu\) be a positive measure on X, F a Banach space, \(F'\) its strong dual, and g any function in \(L_{F'}^{q}\) . Then, the linear form on \(L_{F}^{p}\) , deduced from the linear form \(f \mapsto \int \langle f, g \rangle d\mu\) on \(L_{F}^{p}\) by passage to the quotient, is continuous and has norm \(N_{q}(g)\) .

\subsection{\texorpdfstring{Application: relations between the spaces \(\mathcal{L}_{F}^{p}\) ( \(1 \leqslant p \leqslant +\infty\) )}{Application: relations between the spaces L sub F to the p ( 1 <= p <= +infinity )}}

PROPOSITION 4. --- Let f be a measurable function with values in a Banach space F; the set I of numbers p, such that \(1 \leqslant p \leqslant +\infty\) and \(\mathrm{N}_{p}(\mathbf{f})\) is finite, is either empty or is an interval of \(\overline{R}\) . If I is nonempty, the restriction to I of the mapping \(p \mapsto \mathrm{N}_{p}(\mathbf{f})\) is continuous; moreover, if f is not negligible, \(\log \mathrm{N}_{p}(\mathbf{f})\) is a convex function of 1/p on \(\overline{I}\) .

We already know (Ch. I, No.~3, Prop. 5) that the set J of finite numbers \(p \geqslant 1\) such that \(\mathrm{N}_{p}(\mathbf{f}) < +\infty\) is either empty or is an interval, and that \(\log \mathrm{N}_{p}(\mathbf{f})\) is a convex function of 1/p on J (when f is not negligible); this of course implies the continuity of \(p \mapsto \mathrm{N}_{p}(\mathbf{f})\) on J.

If J is empty then either \(I = \varnothing\) or \(I = \{+\infty\}\) , and the proposition is obvious in this case; assume henceforth that J is nonempty. The proposition is also obvious if f is negligible; assume henceforth that f is not negligible. If \(s \in J\) then, for every finite number p \textgreater{} s, \(|f|^{p} = |f|^{s}|f|^{p-s}\) , and the inequality of the mean shows that

\[
\mathrm{N} _ {p} (\mathbf {f}) \leqslant \left(\mathrm{N} _ {s} (\mathbf {f})\right) ^ {s / p} \left(\mathrm{N} _ {\infty} (\mathbf {f})\right) ^ {(p - s) / p}.\tag{12}
\]

Letting \(p\) tend to \(+\infty\) , it follows that

\[
\limsup _ {p \to + \infty} N _ {p} (\mathbf {f}) \leqslant N _ {\infty} (\mathbf {f})  .\tag{13}
\]

This proves that if \(+\infty \in I\) then \(J\) contains arbitrarily large numbers; thus \(I\) is indeed an interval of \(\overline{\mathbf{R}}\) , and \(\overline{I} = \overline{J}\) . The proposition will be proved if we show that \(p \mapsto N sub p(\mathbf{f})\) is continuous on \(\overline{J}\) , and it suffices to establish continuity at the end-points of \(J\) . We can suppose, moreover, that \(J\) does not reduce to a point. Let \(r\) and \(s\) be the left and right end-points of \(J\) ( \(r < s \leqslant +\infty\) ). Let \(A\) be the (measurable) set of \(x \in X\) such that \(|\mathbf{f}(x)| \geqslant 1\) ; then

\[
\int | \mathbf {f} | ^ {p} d | \mu | = \int | \mathbf {f} | ^ {p} \varphi_ {\mathrm{A}} d | \mu | + \int | \mathbf {f} | ^ {p} \varphi_ {\mathbf {C A}} d | \mu |.
\]

As \(p \in \mathbf{J}\) tends to \(r\) , \(|\mathbf{f}|^p\varphi_{\mathbf{A}}\) tends to \(|\mathbf{f}|^r\varphi_{\mathbf{A}}\) while decreasing, and \(|\mathbf{f}|^p\varphi_{\mathbf{C}_{\mathbf{A}}}\) tends to \(|\mathbf{f}|^r\varphi_{\mathbf{C}_{\mathbf{A}}}\) while increasing. Thus \(\int |\mathbf{f}|^p\varphi_{\mathbf{C}_{\mathbf{A}}}d|\mu|\) tends to \(\int^{*}|\mathbf{f}|^{r}\varphi_{\mathbf{C}_{\mathbf{A}}}d|\mu|\) (§1, No.~3, Th. 3). On the other hand, \(|\mathbf{f}|^p\varphi_{\mathbf{A}}\) is integrable for \(p \in \mathbf{J}\) , and \(\int |\mathbf{f}|^p\varphi_{\mathbf{A}}d|\mu|\) tends to \(\int |\mathbf{f}|^r\varphi_{\mathbf{A}}d|\mu|\) (§4, No.~3, Prop. 4). Therefore \(\int |\mathbf{f}|^p d|\mu|\) tends to \(\int^{*}|\mathbf{f}|^{r}d|\mu|\) , which proves the continuity of \(p \mapsto \mathrm{N}_p(\mathbf{f})\) at \(r\) .

The same reasoning may be applied at the point \(s\) if \(s < +\infty\) . Finally, suppose that \(s = +\infty\) . In view of (13), it suffices to prove that

\[
\liminf _ {p \to + \infty} N _ {p} (\mathbf {f}) \geqslant N _ {\infty} (\mathbf {f}).
\]

Now, let \(a\) be a number such that \(0 < a < \mathrm{N}_{\infty}(\mathbf{f})\) . Since, by hypothesis, there exist finite values of \(p\) such that \(\mathrm{N}_p(\mathbf{f}) < +\infty\) , the set \(A\) of \(x \in X\) such that \(|\mathbf{f}(x)| \geqslant a\) , which is measurable and non-negligible, is integrable by virtue of the inequality \(\varphi_A \leqslant (|\mathbf{f}|/a)^p\) ; moreover, we infer from this inequality that \(\mathrm{N}_p(\mathbf{f}) \geqslant a \cdot (|\mu|(\mathrm{A}))^{1/p}\) ; letting \(p\) tend to \(+\infty\) , it follows that \(\liminf_{p \to +\infty} \mathrm{N}_p(\mathbf{f}) \geqslant a\) , which completes the proof.

COROLLARY. --- If r, s, p are three numbers such that

\[
1 \leqslant r <   p <   s \leqslant + \infty ,
\]

then the intersection \(\mathcal{L}_{\mathrm{F}}^{r} \cap \mathcal{L}_{\mathrm{F}}^{s}\) is contained in \(\mathcal{L}_{\mathrm{F}}^{p}\) .

Note that in general the topologies induced on the intersection \(\mathcal{L}_{\mathrm{F}}^{r} \cap \mathcal{L}_{\mathrm{F}}^{s}\) by the topologies of the \(\mathcal{L}_{\mathrm{F}}^{p}\) ( \(r < p < s\) ) are distinct. If no further hypothesis is made on \(\mu\) , the topologies induced on \(\mathcal{L}_{\mathrm{F}}^{r} \cap \mathcal{L}_{\mathrm{F}}^{s}\) by those of \(\mathcal{L}_{\mathrm{F}}^{r}\) and \(\mathcal{L}_{\mathrm{F}}^{s}\) are in general not comparable (in other words, the ratio \(\mathrm{N}_{r}(\mathbf{f})/\mathrm{N}_{s}(\mathbf{f})\) can take arbitrarily large values and arbitrarily small values in \(L_{F}^{r} \cap L_{F}^{s}\) ; cf.~Exer. 8).

Prop. 4 may be sharpened when \(\mu\) is a bounded measure:

PROPOSITION 5. --- Let \(\mu\) be a bounded measure, and let \(\mathbf{f}\) be a \(\mu\) -measurable function with values in a Banach space \(\mathbf{F}\) . The set I of numbers \(p\) such that \(1 \leqslant p \leqslant +\infty\) and \(\mathrm{N}_p(\mathbf{f})\) is finite, is either empty or is an interval with left end-point \(p = 1\) and containing this point; moreover, \((|\mu|(\mathrm{X}))^{-1/p} \mathrm{N}_p(\mathbf{f})\) is an increasing function of \(p\) on I.

This is an immediate consequence of Prop. 4 above and of the Cor. of Prop. 4 of Ch. I, No.~3.

COROLLARY. --- If the measure \(\mu\) is bounded, the relation r < s implies \(L_{F}^{s} \subset L_{F}^{r}\) ; moreover, the topology of convergence in mean of order s is finer than the topology of convergence in mean of order r (on \(L_{F}^{s}\) ).

One can show that in general the topology of convergence in mean of order \(s\) is strictly finer than the topology of convergence in mean of order \(r\) (Exer. 8).

PROPOSITION 6. --- Let X be a discrete space, \(\mu\) the measure on X defined by placing a mass \(+1\) at each point of X. If \(\mathbf{f}\) is a mapping of X into the Banach space F, the set I of numbers \(p\) such that \(1 \leqslant p \leqslant +\infty\) and \(\mathrm{N}_p(\mathbf{f})\) is finite, is either empty or is an interval with right end-point \(+\infty\) and containing this point; moreover, \(\mathrm{N}_p(\mathbf{f})\) is a decreasing function of \(p\) on I.

For, \(\mu^{*}(|\mathbf{f}|) = \sum_{x\in X}|\mathbf{f}(x)|\) for every function \(\mathbf{f}\) ( \(\S 1\) , No.~1, Example), and \(\mathrm{N}_{\infty}(\mathbf{f}) = \| \mathbf{f}\| = \sup_{x\in X}|\mathbf{f}(x)|\) ; if there exists a number \(\alpha >0\) such that \(|\mathbf{f}(x)|\geqslant \alpha\) for infinitely many values of \(x\in X\) , then \(\mathrm{N}_p(\mathbf{f}) = +\infty\) for every finite \(p\) ; in the contrary case, there exists an \(x_0\in X\) such that \(|\mathbf{f}(x_0)| = \| \mathbf{f}\|\) , whence

\[
\mathrm{N} _ {\infty} (\mathbf {f}) = | \mathbf {f} (x _ {0}) | \leqslant \mathrm{N} _ {p} (\mathbf {f})
\]

for every finite p.~Since the function \(\log\mathrm{N}_{p}(\mathbf{f})\) is convex with respect to 1/p and takes its smallest value at the point \(+\infty\) , it is necessarily a decreasing function of p on I (FRV, I, §4, No.~3, Prop. 5), which completes the proof.

COROLLARY. --- If X is discrete and the measure \(\mu\) is defined by a mass +1 at each point of X, then the relation r < s implies \(L_{F}^{r} \subset L_{F}^{s}\) ; moreover, the topology of convergence in mean of order r is finer than the topology of convergence in mean of order s (on \(L_{F}^{r}\) ).

\section{BARYCENTERS}

\subsection{Definition of barycenters}

Let E be a Hausdorff locally convex space over R, \(E'\) its dual, and \(E'^{*}\) the algebraic dual of \(E'\) , E being canonically identified with a linear subspace of \(E'^{*}\) . Let K be a compact subset of E; the canonical injection of K into E being continuous with compact support, for every measure \(\mu\) on K the integral \(\int x d\mu(x)\) is therefore defined and is an element of \(E'^{*}\) (Ch. III, §3, No.~1). Moreover, on K, the topology induced by the weak topology \(\sigma(E'^{*}, E')\) is identical with the original topology. Finally, if C is the closed convex envelope of K in \(E'^{*}\) equipped with \(\sigma(E'^{*}, E')\) , then \(C \cap E\) is the closed convex envelope of K in E for the original topology (or for the weakened topology \(\sigma(E, E')\) ).

DEFINITION 1. --- Let K be a compact subset of a Hausdorff locally convex space E. For every positive measure \(\mu\) on K of total mass equal to 1, the vector \(\mathbf{b}_{\mu} = \int \mathbf{x} d\mu(\mathbf{x})\) (belonging to \(E'^{*}\) ) is called the barycenter of \(\mu\) .

Example. --- Let \(\mu\) be a discrete measure on K, positive and of total mass 1; it is thus of the form \(\mu = \sum_{i=1}^{n} \lambda_i \varepsilon_{x_i}\) , where \(x_i \in K\) , and the \(\lambda_i\) are real numbers such that \(\lambda_i \geqslant 0\) for all i and \(\sum_{i=1}^{n} \lambda_i = 1\) . Since \(\int x d\varepsilon_y(x) = y\) (Ch. III, §3, No.~1, Example 3), \(b_\mu = \int x d\mu(x) = \sum_{i} \lambda_i x_i\) . In particular, for every \(x \in K\) , x is the barycenter of the measure \(\varepsilon_x\) .

PROPOSITION 1. --- Let E be a Hausdorff locally convex space, K a compact subset of E, and C the closed convex envelope of K in E. The set C then consists of the points of E that are barycenters of at least one positive measure of mass 1 on K.

This is nothing more than Prop. 5 of Ch. III, §3, No.~2 applied to the canonical injection of K into E.

COROLLARY. --- If the closed convex envelope C of K in E is compact, then the barycenter of every positive measure of total mass 1 on K belongs to E.

For, C is then also the closed convex envelope of K in E'* equipped with the weak topology \(\sigma (\mathrm{E}^{\prime *},\mathrm{E}^{\prime})\) , and it suffices to apply, to the canonical injection of K into E, Prop. 4 of Ch. III, §3, No.~2.

Remark. --- The Cor. of Prop. 1 is applicable in particular when K is convex or when E is quasi-complete.

PROPOSITION 2. --- Let K be a compact convex subset of a Hausdorff locally convex space E, \(\mu\) a positive measure of total mass 1 on K, \(b_{\mu}\) its barycenter. For every convex numerical function \(f \geqslant 0\) that is lower semicontinuous on K,

\[
f (\mathbf {b} _ {\mu}) \leqslant \int^ {*} f d \mu .
\]

It is known (TVS, II, §5, No.~4, Prop. 5) that \(f\) is the upper envelope of a family of restrictions to \(K\) of continuous affine linear functions \(h_{\alpha} : \mathbf{x} \mapsto c_{\alpha} + \langle \mathbf{x}, \mathbf{z}_{\alpha}' \rangle\) . Then

\[
\int h _ {\alpha} (\mathbf {x}) d \mu (\mathbf {x}) \leqslant \int^ {*} f (\mathbf {x}) d \mu (\mathbf {x})
\]

for every \(\alpha\) ; now, \(\int h_{\alpha}(\mathbf{x}) d\mu (\mathbf{x}) = c_{\alpha} + \int \langle \mathbf{x}, \mathbf{z}_{\alpha}' \rangle d\mu (\mathbf{x})\) since \(\mu\) has total mass 1; but \(\int \langle \mathbf{x}, \mathbf{z}_{\alpha}' \rangle d\mu (\mathbf{x}) = \langle \mathbf{b}_{\mu}, \mathbf{z}_{\alpha}' \rangle\) by the definition of barycenter. Therefore \(\sup \int h_{\alpha}(\mathbf{x}) d\mu (\mathbf{x}) = \sup h_{\alpha}(\mathbf{b}_{\mu}) = f(\mathbf{b}_{\mu})\) , whence the conclusion.

When \(\mu\) is a discrete positive measure on K of total mass 1, Prop. 2 yields anew the inequality that defines the convex functions on K.

COROLLARY. --- For every convex numerical function \(g\) on \(K\) that is bounded and lower semi-continuous, \(g(\mathbf{b}_{\mu}) \leqslant \int g d\mu\) .

It suffices to note that \(\inf_{\mathbf{x} \in \mathrm{K}} g(\mathbf{x}) = a\) is finite and apply Prop. 2 to \(g - a\) .
